\documentclass[aps,preprint]{revtex4}%
\usepackage{amsfonts}
\usepackage{amsmath}
\usepackage{amssymb}
\usepackage{graphicx}%
\setcounter{MaxMatrixCols}{30}
%TCIDATA{OutputFilter=latex2.dll}
%TCIDATA{Version=5.50.0.2960}
%TCIDATA{CSTFile=revtex4.cst}
%TCIDATA{Created=Friday, May 16, 2003 15:55:00}
%TCIDATA{LastRevised=Sunday, March 14, 2010 10:56:21}
%TCIDATA{<META NAME="GraphicsSave" CONTENT="32">}
%TCIDATA{<META NAME="SaveForMode" CONTENT="1">}
%TCIDATA{BibliographyScheme=Manual}
%TCIDATA{<META NAME="DocumentShell" CONTENT="Articles\SW\REVTeX 4">}
%TCIDATA{Language=American English}
%BeginMSIPreambleData
\providecommand{\U}[1]{\protect\rule{.1in}{.1in}}
%EndMSIPreambleData
\newtheorem{theorem}{Theorem}
\newtheorem{acknowledgement}[theorem]{Acknowledgement}
\newtheorem{algorithm}[theorem]{Algorithm}
\newtheorem{axiom}[theorem]{Axiom}
\newtheorem{claim}[theorem]{Claim}
\newtheorem{conclusion}[theorem]{Conclusion}
\newtheorem{condition}[theorem]{Condition}
\newtheorem{conjecture}[theorem]{Conjecture}
\newtheorem{corollary}[theorem]{Corollary}
\newtheorem{criterion}[theorem]{Criterion}
\newtheorem{definition}[theorem]{Definition}
\newtheorem{example}[theorem]{Example}
\newtheorem{exercise}[theorem]{Exercise}
\newtheorem{lemma}[theorem]{Lemma}
\newtheorem{notation}[theorem]{Notation}
\newtheorem{problem}[theorem]{Problem}
\newtheorem{proposition}[theorem]{Proposition}
\newtheorem{remark}[theorem]{Remark}
\newtheorem{solution}[theorem]{Solution}
\newtheorem{summary}[theorem]{Summary}
\newenvironment{proof}[1][Proof]{\noindent\textbf{#1.} }{\ \rule{0.5em}{0.5em}}
\begin{document}
\preprint{ }
\title[Coh. State Prop.]{A Coherent State View of Propagators}
\author{Brian Weiner}
\affiliation{Department of Physics, Pennsylvania State University, DuBois PA 15801}
\keywords{}
\begin{abstract}


\end{abstract}
\date{06/27/2003}
\startpage{1}
\maketitle
\tableofcontents


\section{Introduction\label{intro}}

Isolated quantum systems do not interact with their environment and are thus
unobservable. In order to determine their properties they must interact with
external probes such as electromagnetic fields of various frequencies. The
radiation field and the system under study, which we will consider to be a
molecular system, can be considered as a joint system that evolves under a
Hamiltonian that describes the behavior of both radiation and molecules.

The effect of the radiation on the molecular subsystem can be monitored by
changes in the molecules density and reduced density matrices. If the
radiation is weak the behaviour of the interaction is well described by the
linear response of the molecular system to the radiation, i.e is linearly
proportional to the radiation strength, while if the radiation is intense the
molecular system responds in a nonlinear way. The interaction of the
electromagnetic field with molecules can be approximated in a standard semi
classical fashion by averaging over the photonic degrees of freedom to produce
effective one electron observables. This leads to a reduced Hamiltonian for
the molecular subsystem that includes perturbing effective time dependent one
body potentials which describe the interaction with the radiation field. In
the case of weak intensity radiation, where the linear response regime is
sufficient to describe the changes in the molecular system, one needs to only
consider responses that are first order in these observables.

The temporal behavior of a state of physical system is completly described by
the evolution of its full and reduced density operators which is determined by
their associated propagators. During the 1970's and 80's a significant effort
was focused by the Yngve Ohrn group at the Quntum Theory Project, University
of Florida on the analysis and implementation of propagator approximations. In
particular the one-electron and polarization propagators were studied by
effectively using the properties of super operators acting on operator
manifolds associated with a given reference state to generate various orders
of approximation \cite{Ohrn, Weiner, Ortiz, Kurtz, Jensen etc}. In this
article we will concentrate on the polarization propagator, which describes
the response of the first order reduced density operator to electromagnetic
radiation. The poles of the polarization propagator give the excitation
energies between stationary states of the molecular system and the
corresponding residues the transition probabilities of these excitations
caused by the electromagnetic perturbation. If the initial molecular reference
state is a stationary state then the poles are the excitation energies and the
residues the transition probabilities from that state. In order for these
approximations to this propagator to be $N$-representable i.e. associated with
some approximate $N$-particle state, the adjoints of the excitation operators
that diagonalize this propagator must annihilate the reference state. However
this in general does not occur and one says that the approximation is
inconsistent. In this article we will consider propagators based on one
particle operators.

If the external radiation fields are of low intensity and the molecular system
is in a stationary state, it is sufficient to consider the linear response of
the molecular state to the applied fields. This can be characterized by a
linear response function that can be calculated by using first order
perturbation theory. Noting that changes in the expectation values of one
particle observables are determined by changes in the one particle reduced
density matrix, one can describe the effects of the radiation field,
(approximated as an effective one particle interaction), on the expectation
values of any one particle observable in terms of retarded two particle propagators.

Two particle propagators based on a given molecular initial state can be
expressed as a second order functional derivative of the propagator, that
contains interaction with arbitrary one particle external fields, with respect
to variations of these fields evaluated at an initial and a final time. These
times are chosen to bracket the time during which the molecular system is
actually interacting, before the initial time and after the final time there
is no interaction and the system is evolving "freely".

In general the quantum propagator of a N particle molecular system
characterizes the linear response of the system to external probes constituted
of fields that interact through M-body forces, where $M\leq N$, as these are
all described by changes in the N-particle density matrix. Thus if one has
explicit analytic expressions for the full propagator of the molecular system
between arbitrary initial and final states one can deduce explicit expressions
for all reduced propagators. If one is interested only to the response of the
system to one body external fields it is sufficient to consider two particle
propagators. This is true for both exact and approximate propagators. As is
well known exact solutions for quantum evolution operators have only been
determined for simple model systems, and are not known for multi electron
molecular systems that involve coulombic interactions between electrons. Thus
for practical applications the goal is to develop systematic, cost effective,
theoretically well founded and non arbitrary approximation schemes, that lead
to accurate enough results to benefit experimentalists.

During the 1970's and 80's we developed approximation methods for two particle
propagators based on the properties of Antisymmetrized Geminal Power (AGP)
states formed from a geminal that was parameterized by the complex variables
$\mathbf{\xi}$. These properties were based on:

\begin{description}
\item[(a)] Excited states,$\left\{  \left\vert \Psi_{n}\right\rangle \right\}
$, of the molecular system being well approximated by single particle
excitations $\left\{  Q_{n}^{\dagger}\right\}  $ from an energy optimized AGP
state i.e.%
\begin{equation}
\left\vert \Psi_{n}\right\rangle =Q_{n}^{\dagger}\left\vert g^{\frac{N}{2}%
}\right\rangle =\sum_{1\leq j,k\leq r}Q_{njk}a_{j}^{\dagger}a_{k}\left\vert
g^{\frac{N}{2}}\right\rangle
\end{equation}


\item[(b)] A vacuum condition being satisfied%
\begin{equation}
Q_{n}\left\vert g^{\frac{N}{2}}\right\rangle =0\label{Killer}%
\end{equation}


\item[(c)] Off diagonal blocks of the energy Hessian
\begin{equation}
D^{2}E\left(  \xi_{0},\bar{\xi}_{0}\right)  =\left(
\begin{array}
[c]{cc}%
\frac{\partial^{2}E\left(  \xi_{0},\bar{\xi}_{0}\right)  }{\partial\xi
\partial\bar{\xi}} & \frac{\partial^{2}E\left(  \xi_{0},\bar{\xi}_{0}\right)
}{\partial\xi\partial\xi}\\
\frac{\partial^{2}E\left(  \xi_{0},\bar{\xi}_{0}\right)  }{\partial\bar{\xi
}\partial\bar{\xi}} & \frac{\partial^{2}E\left(  \xi_{0},\bar{\xi}_{0}\right)
}{\partial\bar{\xi}\partial\bar{\xi}}%
\end{array}
\right)  \label{Hess}%
\end{equation}
of the AGP energy
\begin{equation}
E\left(  \xi,\bar{\xi}\right)  =\frac{\left\langle \left.  g\left(
\xi\right)  ^{\frac{N}{2}}\right\vert Hg\left(  \xi\right)  ^{\frac{N}{2}%
}\right\rangle }{\left\langle \left.  g\left(  \xi\right)  ^{\frac{N}{2}%
}\right\vert g\left(  \xi\right)  ^{\frac{N}{2}}\right\rangle },\label{Energy}%
\end{equation}
at stationary points $\xi_{0}$
\begin{equation}
DE\left(  \xi_{0},\bar{\xi}_{0}\right)  =0
\end{equation}
being small in norm i.e.%
\begin{equation}
\left\Vert \frac{\partial^{2}E\left(  \xi_{0},\bar{\xi}_{0}\right)  }%
{\partial\xi\partial\xi}\right\Vert <\epsilon;\left\Vert \frac{\partial
^{2}E\left(  \xi_{0},\bar{\xi}_{0}\right)  }{\partial\bar{\xi}\partial\bar
{\xi}}\right\Vert <\epsilon.
\end{equation}

\end{description}

It should be noted that the approximation that was developed includes the
perturbing effects of the number operators $\left\{  \left.  a_{j}^{\dagger
}a_{j}\right\vert 1\leq j\leq r\right\}  $, in contrast to the Random Phase
Approximation (RPA) based on Independent Particle States (IPS's). The number
operators describe the response of the energy to canonical coeffient variation
of the geminal, while the operators $\left\{  \left.  a_{j}^{\dagger}%
a_{k}\right\vert 1\leq j\neq k\leq r\right\}  $ describe the response to
orbital variations. In the case of the exact propagator, for which $\left\{
Q_{n}^{\dagger}\right\}  $ are in general multiparticle operators, all of the
properties (a), (b) and (c) are satisfied with the off diagonal Hessian block
being identically zero.

The numerical results obtained from the application of this approximation
scheme agreed well with experiment. 

\section{AGP State}

The reference state we will be focusing on is an AGP\ state given by
\begin{equation}
\left\vert g^{\frac{N}{2}}\right\rangle =\sum_{1\leq j_{1}<j_{2}%
<\cdots<j_{\frac{N}{2}}\leq s}c_{j_{1}}\cdots c_{j_{\frac{N}{2}}}\left\vert
\eta_{j_{1}}\eta_{j_{1}+s}\cdots\eta_{j_{N}\,}\eta_{j_{N}+s}\right\rangle
\label{canexp}%
\end{equation}
where $\left\{  \left.  \left\vert \eta_{j}\right\rangle \right\vert 1\leq
j\leq r\right\}  $ is a complete orthonormal basis for $\mathcal{H}^{1}$ and
\begin{equation}
\left\vert g\right\rangle =\sum_{1\leq j\leq s}c_{j}\left\vert \eta_{j}%
\eta_{j+s}\right\rangle =\sum_{1\leq j\leq s}n_{j}^{\frac{1}{2}}\left\vert
\eta_{j}\eta_{j+s}\right\rangle .
\end{equation}
This state can also be expressed as
\begin{equation}
\left\vert g^{\frac{N}{2}}\right\rangle =P_{\frac{N}{2}}\exp\left\{
\sum_{1\leq j\leq s}\zeta_{j}b_{j}^{\dagger}b_{j+s}^{\dagger}\right\}
\left\vert \phi\right\rangle
\end{equation}
where $P_{\frac{N}{2}}$ is the orthogonal projector onto $\mathcal{H}^{N}$ and
$\left\vert \phi\right\rangle $ is the vacuum state of the Fock space
$\mathcal{H}_{F}$.

\subsection{Group Factorization\label{factor}}

The isotropy group for an AGP reference state is $USp\left(  \mathcal{H}%
^{1}\right)  $ which is a maximally compact subgroup of $SU\left(
\mathcal{H}^{1}\right)  $ hence producing the coset $SU\left(  \mathcal{H}%
^{1}\right)  /USp\left(  \mathcal{H}^{1}\right)  $ One can express coset
representatives in a factored form as $G_{+1}G_{0}G_{-1}$, where $\left\{
G_{+1},G_{0},G_{-1}\right\}  $ are subgroups of $SL\left(  \mathcal{H}%
^{1}\right)  $ and $G_{-1}$ leaves the reference state invariant. Restricting
attention to the connected components of these subgroups, that contain the
identity element , one can express this factorization in terms of Lie algebras
as
\begin{equation}
e^{N_{+}}e^{N_{0}}e^{N_{-}}%
\end{equation}
where $\left\{  N_{+1},N_{0},N_{-1}\right\}  $ are lie subalgebras of
$sl\left(  \mathcal{H}^{1}\right)  $, Introducing coordinate systems $\left\{
\mathbf{z}_{+},\mathbf{\lambda,z}_{-},\mathbf{\eta}\right\}  $ with respect to
a fixed basis of $\left\{  N_{+1},N_{0},N_{-1},\mathrm{h}\right\}  $ where
$\mathrm{h}$ is the isotropy Lie subalgebra $usp\left(  \mathcal{H}%
^{1}\right)  $. One can parameterize group elements $\mathrm{g}\in SU\left(
\mathcal{H}^{1}\right)  $ as
\begin{equation}
g\left(  \mathbf{z,\eta}\right)  =c\left(  \mathbf{z}\right)  h\left(
\mathbf{\eta}\right)  =c_{+1}\left(  \mathbf{z}_{+}\right)  c_{0}\left(
\mathbf{\lambda}\right)  c_{-1}\left(  \mathbf{z}_{-}\right)  h\left(
\mathbf{\eta}\right)
\end{equation}
the non redundant coordinates $\mathbf{z}\,\equiv\left(  \mathbf{z}%
_{+},\mathbf{\lambda}\right)  $ describe a configuration space that
parameterize the set of $\left\{  \left\vert \mathbf{z}\right\rangle \right\}
$ of $SU\left(  \mathcal{H}^{1}\right)  /USp\left(  \mathcal{H}^{1}\right)
$-coherent states and at the same time elements of the coset $SU\left(
\mathcal{H}^{1}\right)  /USp\left(  \mathcal{H}^{1}\right)  $
\begin{equation}
\mathrm{g}\left\vert g^{\frac{N}{2}}\right\rangle =c_{+1}\left(
\mathbf{z}_{+}\right)  c_{0}\left(  \mathbf{\lambda}\right)  c_{-1}\left(
\mathbf{z}_{-}\right)  h\left(  \mathbf{\eta}\right)  \left\vert g^{\frac
{N}{2}}\right\rangle =c_{+1}\left(  \mathbf{z}_{+}\right)  c_{0}\left(
\mathbf{\lambda}\right)  \left\vert g^{\frac{N}{2}}\right\rangle
\end{equation}
the set%
\begin{equation}
\left\vert \mathbf{z}\right\rangle =c_{+1}\left(  \mathbf{z}_{+}\right)
c_{0}\left(  \mathbf{\lambda}\right)  \left\vert g^{\frac{N}{2}}\right\rangle
\end{equation}
of coherent states form a sub manifold $\mathcal{M}$ (in general nonlinear) of
$\mathcal{H}^{N}$ with a coordinate system $\left\{  \mathbf{z}\right\}  $
This manifold is a representation of the coset space $SU\left(  \mathcal{H}%
^{1}\right)  /USp\left(  \mathcal{H}^{1}\right)  .$

The bases one chooses are
\begin{subequations}
\begin{align}
N_{+1} &  =\operatorname{linspan}\left\{  \xi_{j}a_{k}^{\dagger}a_{j}%
+\sigma\left(  j\right)  \sigma\left(  k\right)  \xi_{k}a_{-j}^{\dagger}%
a_{-k};-s\leq j<k\leq s,\left\vert j\right\vert \neq\left\vert k\right\vert
;\,\sigma\left(  j\right)  =\operatorname{sign}\left(  j\right)  \right\}  \\
N_{0} &  =\operatorname{linspan}\left\{  a_{k}^{\dagger}a_{j}-\frac{1}%
{2}\delta_{kj}I;\quad-s\leq j\leq s,\left\vert j\right\vert =\left\vert
k\right\vert \right\}  \\
N_{-1} &  =\operatorname{linspan}\left\{  \xi_{j}a_{k}^{\dagger}a_{j}%
-\sigma\left(  j\right)  \sigma\left(  k\right)  \xi_{k}a_{-j}^{\dagger}%
a_{-k};-s\leq j<k\leq s,\left\vert j\right\vert \neq\left\vert k\right\vert
;\,\sigma\left(  j\right)  =\operatorname{sign}\left(  j\right)  \right\}  \\
\mathrm{h} &  =\operatorname{linspan}\left\{  i\left(  a_{k}^{\dagger}%
a_{j}-\frac{1}{2}\delta_{kj}I\right)  ;\quad-s\leq j\leq s,\left\vert
j\right\vert =\left\vert k\right\vert \right\}
\end{align}
and the factors in the group decomposition are%
\end{subequations}
\begin{align}
c_{+1}\left(  \mathbf{z}_{+}\right)   &  =\exp\left\{  \sum
\limits_{\substack{-s\leq j<k\leq s\\\left\vert j\right\vert \neq\left\vert
k\right\vert }}z_{+jk}\left(  \xi_{j}a_{k}^{\dagger}a_{j}+\sigma\left(
j\right)  \sigma\left(  k\right)  \xi_{k}a_{-j}^{\dagger}a_{-k}\right)
\right\}  =\exp\left\{  \sum\limits_{_{\substack{-s\leq j<k\leq s\\\left\vert
j\right\vert \neq\left\vert k\right\vert }}}z_{+jk}q_{jk}^{\dagger}\right\}
\\
c_{-1}\left(  \mathbf{z}_{-}\right)   &  =\exp\left\{  -\sum
\limits_{\substack{-s\leq j<k\leq s\\\left\vert j\right\vert \neq\left\vert
k\right\vert }}z_{-jk}\left(  \xi_{j}a_{k}^{\dagger}a_{j}-\sigma\left(
j\right)  \sigma\left(  k\right)  \xi_{k}a_{-j}^{\dagger}a_{-k}\right)
\right\}  =\exp\left\{  \sum\limits_{_{\substack{-s\leq j<k\leq s\\\left\vert
j\right\vert \neq\left\vert k\right\vert }}}z_{-jk}q_{jk}\right\}  \\
c_{0}\left(  \mathbf{\lambda}\right)   &  =\exp\left\{  \sum
\limits_{_{\substack{-s\leq j<k\leq s\\\left\vert j\right\vert =\left\vert
k\right\vert }}}\lambda_{jk}\left(  a_{k}^{\dagger}a_{j}-\frac{1}{2}%
\delta_{kj}I\right)  \right\}  =\exp\left\{  \sum\limits_{_{\substack{-s\leq
j<k\leq s\\\left\vert j\right\vert =\left\vert k\right\vert }}}\lambda
_{jk}u_{jk}\right\}  ;\quad\lambda_{jk}=\bar{\lambda}_{kj}\\
h\left(  \mathbf{\eta}\right)   &  =\exp\left\{  \sum
\limits_{_{\substack{-s\leq j<k\leq s\\\left\vert j\right\vert =\left\vert
k\right\vert }}}\eta_{jk}\left(  a_{k}^{\dagger}a_{j}-\frac{1}{2}\delta
_{kj}I\right)  \right\}  =\exp\left\{  \sum\limits_{_{\substack{-s\leq j<k\leq
s\\\left\vert j\right\vert =\left\vert k\right\vert }}}\eta_{jk}%
u_{jk}\right\}  ;\quad\eta_{jk}=-\bar{\eta}_{kj}%
\end{align}


Normalized coherent states are produced by the action of $G$ on normalized
reference states and thus by the action of the coset representatives. Often
however it is convenient to consider unnormalized coherent states by either
dropping or modifying the coset representative $c_{0}\left(  \mathbf{\lambda
}\right)  $ and explicitly considering the norm of the state.

and is induced by the K\"{a}hler potential
\begin{equation}
F\left(  \mathbf{z}^{\prime},\mathbf{z}\right)  =\ln\left\langle
\mathbf{z}^{\prime}|\mathbf{z}\right\rangle
\end{equation}
and defines a generalized Poisson Bracket%
\begin{equation}
\left\{  f,g\right\}  =\omega\left(  X_{f},X_{g}\right)
\end{equation}
The generalized classical equations of motion defined on this phase space can
then be expressed as
\begin{align}
\frac{i}{\hslash}\frac{dz_{j}}{dt} &  =\left\{  z_{j},E\right\}  \nonumber\\
-\frac{i}{\hslash}\frac{d\bar{z}_{j}}{dt} &  =\left\{  \bar{z}_{j},E\right\}
\label{EOM}%
\end{align}
which can be written in matrix form as
\begin{equation}
\left(
\begin{array}
[c]{c}%
\frac{d\mathbf{\bar{z}}}{dt}\\
\frac{d\mathbf{z}}{dt}%
\end{array}
\right)  =\left(
\begin{array}
[c]{cc}%
\mathbf{0} & \frac{i}{\hslash}\mathbf{\bar{C}}^{-1}\\
-\frac{i}{\hslash}\mathbf{C}^{-1} & \mathbf{0}%
\end{array}
\right)  \left(
\begin{array}
[c]{c}%
\frac{\partial E}{\partial\mathbf{\bar{z}}}\\
\frac{\partial E}{\partial\mathbf{z}}%
\end{array}
\right)
\end{equation}
The metric matrix $\mathbf{C}=\left.  \frac{\partial^{2}\ln\left\langle
\mathbf{z}^{\prime}|\mathbf{z}\right\rangle }{\partial\bar{z}_{j}\partial
z_{k}}\right\vert _{\mathbf{z=z\prime}}$ is given by
\begin{equation}
\mathbf{C}\left(  \mathbf{z},\mathbf{\bar{z}}\right)  =\frac{1}{S\left(
\mathbf{z},\mathbf{\bar{z}}\right)  }\left\{  \left.  \frac{\partial
^{2}S\left(  \mathbf{z}^{\prime},\mathbf{z}\right)  }{\partial\mathbf{\bar{z}%
}\partial\mathbf{z}}\right\vert _{\mathbf{z=z\prime}}-\frac{1}{S\left(
\mathbf{z},\mathbf{\bar{z}}\right)  }\left.  \frac{\partial S\left(
\mathbf{z}^{\prime},\mathbf{z}\right)  }{\partial\mathbf{\bar{z}}}\right\vert
_{\mathbf{z=z\prime}}\left.  \frac{\partial S\left(  \mathbf{z}^{\prime
},\mathbf{z}\right)  }{\partial\mathbf{z}}\right\vert _{\mathbf{z=z\prime}%
}\right\}
\end{equation}
where $S\left(  \mathbf{z}^{\prime},\mathbf{z}\right)  =\left\langle
\mathbf{z}^{\prime}|\mathbf{z}\right\rangle $.

The gradient of the energy function $E\left(  \mathbf{z},\mathbf{\bar{z}%
}\right)  $ has the form%
\begin{equation}
DE\left(  \mathbf{z},\mathbf{\bar{z}}\right)  =\left(
\begin{array}
[c]{c}%
\frac{\partial E\left(  \mathbf{z},\mathbf{\bar{z}}\right)  }{\partial
\mathbf{\bar{z}}}\\
\frac{\partial E\left(  \mathbf{z},\mathbf{\bar{z}}\right)  }{\partial
\mathbf{z}}%
\end{array}
\right)
\end{equation}
where
\begin{equation}
\frac{\partial E\left(  \mathbf{z},\mathbf{\bar{z}}\right)  }{\partial
\mathbf{x}}=\frac{1}{S\left(  \mathbf{z},\mathbf{\bar{z}}\right)  }\left\{
\frac{\partial H\left(  \mathbf{z},\mathbf{\bar{z}}\right)  }{\partial
\mathbf{x}}-E\left(  \mathbf{z},\mathbf{\bar{z}}\right)  \frac{\partial
S\left(  \mathbf{z},\mathbf{\bar{z}}\right)  }{\partial\mathbf{x}}\right\}
\end{equation}
for $\mathbf{x=z,\bar{z}}$. At critical points $\left(  \mathbf{z}%
_{c},\mathbf{\bar{z}}_{c}\right)  $ of $E\left(  \mathbf{z},\mathbf{\bar{z}%
}\right)  $
\begin{equation}
\frac{1}{S\left(  \mathbf{z}_{c},\mathbf{\bar{z}}_{c}\right)  }\left\{
\frac{\partial H\left(  \mathbf{z}_{c},\mathbf{\bar{z}}_{c}\right)  }%
{\partial\mathbf{x}}-E\left(  \mathbf{z}_{c},\mathbf{\bar{z}}_{c}\right)
\frac{\partial S\left(  \mathbf{z}_{c},\mathbf{\bar{z}}_{c}\right)  }%
{\partial\mathbf{x}}\right\}  =0
\end{equation}
hence
\begin{equation}
\frac{\partial H\left(  \mathbf{z}_{c},\mathbf{\bar{z}}_{c}\right)  }%
{\partial\mathbf{x}}=E\left(  \mathbf{z}_{c},\mathbf{\bar{z}}_{c}\right)
\frac{\partial S\left(  \mathbf{z}_{c},\mathbf{\bar{z}}_{c}\right)  }%
{\partial\mathbf{x}}%
\end{equation}


The Hessian $D^{2}E\left(  \mathbf{z},\mathbf{\bar{z}}\right)  $ of the energy
at critical points is then made up of blocks of the form
\begin{equation}
\frac{\partial^{2}E\left(  \mathbf{z}_{c},\mathbf{\bar{z}}_{c}\right)
}{\partial\mathbf{x}\partial\mathbf{y}}=\frac{1}{S\left(  \mathbf{z}%
_{c},\mathbf{\bar{z}}_{c}\right)  }\left\{  \frac{\partial^{2}H\left(
\mathbf{z}_{c},\mathbf{\bar{z}}_{c}\right)  }{\partial\mathbf{x}%
\partial\mathbf{y}}-E\left(  \mathbf{z}_{c},\mathbf{\bar{z}}_{c}\right)
\frac{\partial^{2}S\left(  \mathbf{z}_{c},\mathbf{\bar{z}}_{c}\right)
}{\partial\mathbf{x}\partial\mathbf{y}}\right\}
\end{equation}
for $\mathbf{x},\mathbf{y=z},\mathbf{\bar{z}}$.

\section{Linear Response\label{linres}}

\subsection{Exact Linear Response\label{exlinres}}

Exact linear response functions are defined as
\begin{equation}
R_{FA}\left(  t\right)  =-\frac{i}{\hslash}\theta_{+}\left(  t\right)
\left\langle \Psi_{0}\left\vert \left[  F_{H}\left(  t\right)  ,A\left(
0\right)  \right]  \right.  \Psi_{0}\right\rangle
\end{equation}
which describes the response of the expectation value of the observable $A$ to
the perturbation $F,$ where $\left\vert \Psi_{0}\right\rangle $ is an exact
stationary state of an unperturbed Hamiltonian $H_{o}$, and the time
dependence of observable $A$ is given by
\begin{equation}
A\left(  t\right)  =e^{-iH_{o}t}Ae^{iH_{o}t}.
\end{equation}
and
\begin{equation}
F_{H}\left(  t\right)  =e^{-iH_{o}t}F\left(  t\right)  e^{iH_{o}t}%
\end{equation}
If $A$ and $F$ are one-particle operators this response can be formulated in
terms of the response of first order reduced density matrix elements to
external time dependent perturbing fields, $F(t)$, that interact with the
system through one body forces, according to the following:%
\begin{equation}
F(t)=\sum_{k,l}f_{kl}\left(  t\right)  a_{k}^{\dagger}\left(  t\right)
a_{l}\left(  t\right)
\end{equation}
where
\begin{equation}
a_{k}^{\dagger}\left(  t\right)  a_{l}\left(  t\right)  =e^{-iH_{o}t}%
a_{k}^{\dagger}a_{l}e^{iH_{o}t}%
\end{equation}%
\begin{equation}
\rho_{ij}\left(  t\right)  =\rho_{ij}\left(  0\right)  +\frac{i}{\hslash}%
\int_{0}^{t}\left\langle \Psi_{0}\left\vert \left[  F\left(  \tau\right)
,a_{k}^{\dagger}a_{l}\right]  \right.  \Psi_{0}\right\rangle d\tau
\end{equation}%
\begin{equation}
\rho_{ij}\left(  t\right)  =\rho_{ij}\left(  0\right)  +\delta\rho_{ij}\left(
t\right)
\end{equation}%
\begin{equation}
\delta\rho_{ij}\left(  t\right)  =\sum_{k,l}\int_{-\infty}^{+\infty}%
R_{ijkl}\left(  t-\tau\right)  f_{kl}\left(  \tau\right)  d\tau
\end{equation}
where
\begin{equation}
R_{ijkl}\left(  t\right)  =-\frac{i}{\hslash}\left\langle \Psi_{0}\left\vert
\left[  a_{j}^{\dagger}\left(  t\right)  a_{i}\left(  t\right)  ,a_{k}%
^{\dagger}a_{l}\right]  \right.  \Psi_{0}\right\rangle
\end{equation}
$\rho_{ij}$ are first order reduced density matrix elements and $R_{ijkl}%
\left(  t\right)  $ are the one particle linear response functions. The
response function may then be written as
\begin{equation}
R_{ijkl}\left(  t\right)  =-\frac{i}{\hslash}\theta_{+}\left(  t\right)
\left\{  \left\langle \Psi_{0}\left\vert a_{j}^{\dagger}a_{i}e^{-\frac
{i}{\hslash}\left(  H_{0}-E_{0}\right)  t}a_{k}^{\dagger}a_{l}\right.
\Psi_{0}\right\rangle -\left\langle \Psi_{0}\left\vert a_{k}^{\dagger}%
a_{l}e^{-\frac{i}{\hslash}\left(  H_{0}-E_{0}\right)  t}a_{j}^{\dagger}%
a_{i}\right.  \Psi_{0}\right\rangle \right\}
\end{equation}
where $E_{0}$ is the energy of the stationary state $\left\vert \Psi
_{0}\right\rangle $. Taking the Laplace transform of $R_{ijkl}\left(
t\right)  $ we obtain
\begin{equation}
R_{ijkl}\left(  z\right)  =\sum_{n\neq0}\left\{  \frac{\left\langle \Psi
_{0}\left\vert a_{j}^{\dagger}a_{i}\Psi_{n}\right.  \right\rangle \left\langle
\Psi_{n}\left\vert a_{k}^{\dagger}a_{l}\Psi_{0}\right.  \right\rangle
}{z-\left(  E_{n}-E_{0}\right)  }-\frac{\left\langle \Psi_{0}\left\vert
a_{k}^{\dagger}a_{l}\Psi_{n}\right.  \right\rangle \left\langle \Psi
_{n}\left\vert a_{j}^{\dagger}a_{k}\Psi_{0}\right.  \right\rangle }{z+\left(
E_{n}-E_{0}\right)  }\right\}
\end{equation}
The response functions $R_{ijkl}\left(  z\right)  $ form a hermitian matrix
for real $z$, thus can be diagonalized to give
\begin{equation}
\tilde{R}_{KL}\left(  z\right)  =g_{K}\left(  z\right)  \delta_{KL}%
=\sum_{n\neq0}\left\{  \frac{\left\langle \Psi_{0}\left\vert T_{K}\Psi
_{n}\right.  \right\rangle \left\langle \Psi_{n}\left\vert T_{K}^{\dagger}%
\Psi_{0}\right.  \right\rangle }{z-\left(  E_{n}-E_{0}\right)  }%
-\frac{\left\langle \Psi_{0}\left\vert T_{K}\Psi_{n}\right.  \right\rangle
\left\langle \Psi_{n}\left\vert T_{K}^{\dagger}\Psi_{0}\right.  \right\rangle
}{z+\left(  E_{n}-E_{0}\right)  }\right\}
\end{equation}
where
\begin{equation}
T_{K}^{\dagger}=\sum_{ij}T_{Kij}a_{i}^{\dagger}a_{j}%
\end{equation}
The poles of the Laplace transform occur at $\pm\left(  E_{n}-E_{0}\right)  $
and the residues $\left\vert \left\langle \Psi_{n}\left\vert T_{K}^{\dagger
}\Psi_{0}\right.  \right\rangle \right\vert ^{2}$ are transition probabilities
that the operator will cause a transition from the stationary state
$\left\vert \Psi_{0}\right\rangle $ to $\left\vert \Psi_{n}\right\rangle $ the
stationary state. In general the excited states are of the form%
\begin{equation}
\left\vert \Psi_{n}\right\rangle =\sum_{0\leq m\leq N}\sum_{K_{1}\cdots K_{m}%
}c_{K_{1}\cdots K_{m}}^{n}T_{K_{1}}^{\dagger}\cdots T_{K_{m}}^{\dagger
}\left\vert \Psi_{0}\right\rangle
\end{equation}
In this study we focus on excited states that are produced from the reference
state by single photon processes, i.e. low intensity electromagnetic
radiation, which is consistent with the hypotheses of linear response. Thus we
desire that the excited states be well approximated by%
\begin{equation}
\left\vert \Psi_{n}\right\rangle \approx T_{n}^{\dagger}\left\vert \Psi
_{0}\right\rangle
\end{equation}


\subsection{Linear Response in a Coherent State Setting}

In the coherent state setting one can see the implications of these
approximations more clearly, i.e. the approx. are more precisiely defined and
not arbitrary.

The classical Hamiltonian corresponding to the Hamiltonian containing the time
dependent perturbation is
\begin{equation}
E\left(  \mathbf{z}\left(  t\right)  ,\mathbf{\bar{z}}\left(  t\right)
\mathbf{,}t\right)  =\frac{\left\langle \left.  \mathbf{z}\left(  t\right)
\right\vert \left(  H_{0}+F\left(  t\right)  \right)  \mathbf{z}\left(
t\right)  \right\rangle }{\left\langle \left.  \mathbf{z}\left(  t\right)
\right\vert \mathbf{z}\left(  t\right)  \right\rangle }=E_{0}\left(
\mathbf{z},\mathbf{\bar{z}}\right)  +\frac{\left\langle \left.  \mathbf{z}%
\right\vert \left(  t\right)  F\left(  t\right)  \mathbf{z}\left(  t\right)
\right\rangle }{\left\langle \left.  \mathbf{z}\left(  t\right)  \right\vert
\mathbf{z}\left(  t\right)  \right\rangle }%
\end{equation}
One can form a quadratic approximation to this function around a point
$\left(  \mathbf{z}_{0},\mathbf{\bar{z}}_{0}\right)  $ by expanding up to
second order
\begin{align}
E\left(  \mathbf{z}_{0}+\delta\mathbf{z}\left(  t\right)  ,\mathbf{\bar{z}%
}_{0}+\delta\mathbf{\bar{z}}\left(  t\right)  \right)   &  \approx
E_{2}\left[  \mathbf{z}_{0}\right]  \left(  \delta\mathbf{z}\left(  t\right)
,\delta\mathbf{\bar{z}}\left(  t\right)  ,t\right) \\
&  =E_{0}\left(  \mathbf{z}_{0},\mathbf{\bar{z}}_{0},t\right)  +\delta
^{1}E_{0}\left(  \mathbf{z}_{0},\mathbf{\bar{z}}_{0},t\right)  +\delta
^{2}E_{0}\left(  \mathbf{z}_{0},\mathbf{\bar{z}}_{0},t\right)
\end{align}
where
\begin{align}
\delta^{1}E_{0}\left(  \mathbf{z}_{0},\mathbf{\bar{z}}_{0},t\right)   &
=D^{1}E\left(  \mathbf{z}_{0},\mathbf{\bar{z}}_{0},t\right)  \left(
\begin{array}
[c]{c}%
\delta\mathbf{z}\left(  t\right) \\
\delta\mathbf{\bar{z}}\left(  t\right)
\end{array}
\right) \\
\delta^{2}E_{0}\left(  \mathbf{z}_{0},\mathbf{\bar{z}}_{0},t\right)   &
=\frac{1}{2}\left(
\begin{array}
[c]{cc}%
\delta\mathbf{z}\left(  t\right)  & \delta\mathbf{\bar{z}}\left(  t\right)
\end{array}
\right)  D^{2}E\left(  \mathbf{z}_{0},\mathbf{\bar{z}}_{0},t\right)  \left(
\begin{array}
[c]{c}%
\delta\mathbf{z}\left(  t\right) \\
\delta\mathbf{\bar{z}}\left(  t\right)
\end{array}
\right)
\end{align}
The approximate classical Hamiltonian associated with linear response about
the point $\left(  \mathbf{z}_{0},\mathbf{\bar{z}}_{0}\right)  $ is then
constructed by only including the time dependent perturbation in terms up to
first order in $\left(  \delta\mathbf{z},\delta\mathbf{\bar{z}}\right)  $
\begin{align}
&  E_{2}\left[  \mathbf{z}_{0}\right]  \left(  \delta\mathbf{z}\left(
t\right)  ,\delta\mathbf{\bar{z}}\left(  t\right)  ,t\right)  =E_{0}\left(
\mathbf{z}_{0},\mathbf{\bar{z}}_{0},t\right)  +\left(
\begin{array}
[c]{c}%
\frac{\partial E\left(  \mathbf{z}_{0},\mathbf{\bar{z}}_{0},t\right)
}{\partial\mathbf{\bar{z}}}\\
\frac{\partial E\left(  \mathbf{z}_{0},\mathbf{\bar{z}}_{0},t\right)
}{\partial\mathbf{z}}%
\end{array}
\right)  \left(
\begin{array}
[c]{c}%
\delta\mathbf{\bar{z}}\left(  t\right) \\
\delta\mathbf{z}\left(  t\right)
\end{array}
\right) \nonumber\\
&  +\frac{1}{2}\left(
\begin{array}
[c]{cc}%
\delta\mathbf{\bar{z}}\left(  t\right)  & \delta\mathbf{z}\left(  t\right)
\end{array}
\right)  \left(
\begin{array}
[c]{cc}%
\frac{\partial^{2}E\left(  \mathbf{z}_{0},\mathbf{\bar{z}}_{0},t\right)
}{\partial\mathbf{\bar{z}}\partial\mathbf{\bar{z}}} & \frac{\partial
^{2}E\left(  \mathbf{z}_{0},\mathbf{\bar{z}}_{0},t\right)  }{\partial
\mathbf{\bar{z}}\partial\mathbf{z}}\\
\frac{\partial^{2}E\left(  \mathbf{z}_{0},\mathbf{\bar{z}}_{0},t\right)
}{\partial\mathbf{z}\partial\mathbf{\bar{z}}} & \frac{\partial^{2}E\left(
\mathbf{z}_{0},\mathbf{\bar{z}}_{0},t\right)  }{\partial\mathbf{z}%
\partial\mathbf{z}}%
\end{array}
\right)  \left(
\begin{array}
[c]{c}%
\delta\mathbf{\bar{z}}\left(  t\right) \\
\delta\mathbf{z}\left(  t\right)
\end{array}
\right)
\end{align}
using the linear response approximation one obtains an approximate gradient
for local variations as
\begin{equation}
\left(
\begin{array}
[c]{c}%
\frac{\partial E_{2}\left[  \mathbf{z}_{0}\right]  \left(  \delta
\mathbf{z},\delta\mathbf{\bar{z}},t\right)  }{\partial\delta\mathbf{\bar{z}}%
}\\
\frac{\partial E_{2}\left[  \mathbf{z}_{0}\right]  \left(  \delta
\mathbf{z},\delta\mathbf{\bar{z}},t\right)  }{\partial\delta\mathbf{z}}%
\end{array}
\right)  =D^{1}E\left(  \mathbf{z}_{0},\mathbf{\bar{z}}_{0},t\right)
+D^{2}E\left(  \mathbf{z}_{0},\mathbf{\bar{z}}_{0},t\right)  \left(
\begin{array}
[c]{c}%
\delta\mathbf{\bar{z}}\left(  t\right) \\
\delta\mathbf{z}\left(  t\right)
\end{array}
\right)
\end{equation}
which can be used to obtain a linearized EOM in the tangent space at the point
$\left(  \mathbf{z}_{0},\mathbf{\bar{z}}_{0}\right)  $
\begin{equation}
\left(
\begin{array}
[c]{c}%
\frac{d\delta\mathbf{\bar{z}}}{dt}\\
\frac{d\delta\mathbf{z}}{dt}%
\end{array}
\right)  =\left(
\begin{array}
[c]{cc}%
\mathbf{0} & \frac{i}{\hslash}\mathbf{\bar{C}}^{-1}\\
-\frac{i}{\hslash}\mathbf{C}^{-1} & \mathbf{0}%
\end{array}
\right)  \left(  D^{1}E\left(  \mathbf{z}_{0},\mathbf{\bar{z}}_{0},t\right)
+D^{2}E\left(  \mathbf{z}_{0},\mathbf{\bar{z}}_{0},t\right)  \left(
\begin{array}
[c]{c}%
\delta\mathbf{\bar{z}}\left(  t\right) \\
\delta\mathbf{z}\left(  t\right)
\end{array}
\right)  \right)
\end{equation}
If the point $\left(  \mathbf{z}_{0},\mathbf{\bar{z}}_{0}\right)  $ is a
critical point of the unperturbed energy $E\left(  \mathbf{z}_{0}%
,\mathbf{\bar{z}}_{0}\right)  $ then
\begin{align}
\frac{\partial E\left(  \mathbf{z}_{0},\mathbf{\bar{z}}_{0},0\right)
}{\partial\mathbf{x}}  &  =\frac{1}{S\left(  \mathbf{z}_{0},\mathbf{\bar{z}%
}_{0}\right)  }\left\{  \frac{\partial\mathcal{F}\left(  \mathbf{z}%
_{0},\mathbf{\bar{z}}_{0},0\right)  }{\partial\mathbf{x}}-\frac{\left\langle
\left.  \mathbf{z}_{0}\right\vert F\left(  0\right)  \mathbf{z}_{0}%
\right\rangle }{\left\langle \left.  \mathbf{z}_{0}\right\vert \mathbf{z}%
_{0}\right\rangle }\frac{\partial S\left(  \mathbf{z}_{0},\mathbf{\bar{z}}%
_{0},t\right)  }{\partial\mathbf{x}}\right\} \\
&  =\frac{1}{S\left(  \mathbf{z}_{0},\mathbf{\bar{z}}_{0}\right)  }%
\frac{\partial F\left(  \mathbf{z}_{0},\mathbf{\bar{z}}_{0},0\right)
}{\partial\mathbf{x}}%
\end{align}
where $\mathcal{F}\left(  \mathbf{z}_{0},\mathbf{\bar{z}}_{0},t\right)
=\left\langle \left.  \mathbf{z}_{0}\right\vert F\left(  t\right)
\mathbf{z}_{0}\right\rangle ,$ $\mathbf{z}_{0}=\mathbf{z}\left(  0\right)  $
$\lceil$as $\mathbf{z}\left(  t\right)  =\mathbf{z}_{0}+\delta\mathbf{z}%
\left(  t\right)  \rfloor$ and at $t=0$ the external perturbation $F\left(
0\right)  =0$. The Hessian at this point is
\begin{equation}
\frac{\partial^{2}E\left(  \mathbf{z}_{0},\mathbf{\bar{z}}_{0},0\right)
}{\partial\mathbf{x}\partial\mathbf{y}}=\frac{1}{S\left(  \mathbf{z}%
_{0},\mathbf{\bar{z}}_{0},0\right)  }\left\{  \frac{\partial^{2}H_{0}\left(
\mathbf{z}_{0},\mathbf{\bar{z}}_{0}\right)  }{\partial\mathbf{x}%
\partial\mathbf{y}}-E\left(  \mathbf{z}_{0},\mathbf{\bar{z}}_{0},0\right)
\frac{\partial^{2}S\left(  \mathbf{z}_{0},\mathbf{\bar{z}}_{0},0\right)
}{\partial\mathbf{x}\partial\mathbf{y}}\right\}
\end{equation}
where $\mathbf{x},\mathbf{y=z},\mathbf{\bar{z}.}$ The linear response $\left(
\delta\mathbf{z},\delta\mathbf{\bar{z}}\right)  $ of the state around a
critical point is then determined by the equation
\begin{align}
\left(
\begin{array}
[c]{cc}%
\mathbf{0} & \frac{i}{\hslash}\mathbf{C}\\
-\frac{i}{\hslash}\mathbf{\bar{C}} & \mathbf{0}%
\end{array}
\right)  \left(
\begin{array}
[c]{c}%
\frac{d\delta\mathbf{\bar{z}}}{dt}\\
\frac{d\delta\mathbf{z}}{dt}%
\end{array}
\right)   &  =\left(
\begin{array}
[c]{c}%
\frac{\partial E_{2}\left[  \mathbf{z}_{0}\right]  \left(  \delta
\mathbf{z},\delta\mathbf{\bar{z}},t\right)  }{\partial\delta\mathbf{\bar{z}}%
}\\
\frac{\partial E_{2}\left[  \mathbf{z}_{0}\right]  \left(  \delta
\mathbf{z},\delta\mathbf{\bar{z}},t\right)  }{\partial\delta\mathbf{z}}%
\end{array}
\right) \nonumber\\
&  =D^{1}F\left(  \mathbf{z}_{0},\mathbf{\bar{z}}_{0},t\right)  +D^{2}E\left(
\mathbf{z}_{0},\mathbf{\bar{z}}_{0},t\right)  \left(
\begin{array}
[c]{c}%
\delta\mathbf{\bar{z}}\left(  t\right) \\
\delta\mathbf{z}\left(  t\right)
\end{array}
\right)
\end{align}
which gives that
\begin{equation}
\left\{  \left(
\begin{array}
[c]{cc}%
\mathbf{0} & \frac{i}{\hslash}\mathbf{C}\frac{d}{dt}\\
-\frac{i}{\hslash}\mathbf{\bar{C}}\frac{d}{dt} & \mathbf{0}%
\end{array}
\right)  -D^{2}E\left(  \mathbf{z}_{0},\mathbf{\bar{z}}_{0},t\right)
\right\}  \left(
\begin{array}
[c]{c}%
\delta\mathbf{\bar{z}}\left(  t\right) \\
\delta\mathbf{z}\left(  t\right)
\end{array}
\right)  =D^{1}F\left(  \mathbf{z}_{0},\mathbf{\bar{z}}_{0},t\right)
\end{equation}
and
\begin{equation}
\left(
\begin{array}
[c]{c}%
\delta\mathbf{\bar{z}}\left(  t\right) \\
\delta\mathbf{z}\left(  t\right)
\end{array}
\right)  =\int_{-\infty}^{+\infty}R\left(  t-\tau\right)  D^{1}F\left(
\mathbf{z}_{0},\mathbf{\bar{z}}_{0},t\right)  d\tau
\end{equation}
where the response function is defined as
\begin{equation}
R\left(  t\right)  =\theta_{+}\left(  t\right)  \left\{  \left(
\begin{array}
[c]{cc}%
\mathbf{0} & \frac{i}{\hslash}\mathbf{C}\frac{d}{dt}\\
-\frac{i}{\hslash}\mathbf{\bar{C}}\frac{d}{dt} & \mathbf{0}%
\end{array}
\right)  -D^{2}E\left(  \mathbf{z}_{0},\mathbf{\bar{z}}_{0},t\right)
\right\}  ^{-1}%
\end{equation}
Taking the Laplace transform of $R\left(  t\right)  $ we obtain
\begin{equation}
\left(
\begin{array}
[c]{c}%
\delta\mathbf{\bar{z}}\left(  \xi\right) \\
\delta\mathbf{z}\left(  \xi\right)
\end{array}
\right)  =R\left(  \xi\right)  D^{1}F\left(  \mathbf{z}_{0},\mathbf{\bar{z}%
}_{0},\xi\right)
\end{equation}
where
\begin{equation}
R\left(  \xi\right)  =\left\{  \left(
\begin{array}
[c]{cc}%
\mathbf{0} & -\xi\mathbf{C}\\
\bar{\xi}\mathbf{\bar{C}} & \mathbf{0}%
\end{array}
\right)  -D^{2}E\left(  \mathbf{z}_{0},\mathbf{\bar{z}}_{0}\right)  \right\}
^{-1}%
\end{equation}
The preceding response function $R\left(  t\right)  $ describes the response
of observables of a system described by $H_{0}$ to external perturbations
$F\left(  t\right)  $ that interact through 1,2,....,N body forces. The
response of the observables are described by the response of the density and
reduced density matrices
\begin{equation}
\left\{  D^{N}\left(  \mathbf{z}_{0},\mathbf{\bar{z}}_{0}\right)
,\ldots,D^{1}\left(  \mathbf{z}_{0},\mathbf{\bar{z}}_{0}\right)  \right\}
\end{equation}
The linear response of the matrix elements $D^{N}\left(  \mathbf{z}%
_{0},\mathbf{\bar{z}}_{0}\right)  _{KL}$ of the N-particle density operator
$\left\vert \mathbf{z}_{0}\right\rangle \left\langle \mathbf{z}_{0}\right\vert
$ of the coherent state $\left\vert \mathbf{z}_{0}\right\rangle $ to the
external perturbation are given by
\begin{align}
\delta^{1}\left\langle \Phi_{K}\right.  \left\vert \mathbf{z}_{0}\right\rangle
\left\langle \mathbf{z}_{0}\right\vert \left.  \Phi_{L}\right\rangle \left(
t\right)   &  =\delta^{1}\left\langle \mathbf{z}_{0}\right\vert \left.
\Phi_{L}\right\rangle \left\langle \Phi_{K}\right.  \left\vert \mathbf{z}%
_{0}\right\rangle \left(  t\right)  \equiv\delta^{1}\Phi_{KL}\left(
\mathbf{z}_{0},\mathbf{\bar{z}}_{0}\right) \nonumber\\
&  =\left(
\begin{array}
[c]{cc}%
\frac{\partial\Phi_{KL}\left(  \mathbf{z}_{0},\mathbf{\bar{z}}_{0}\right)
}{\partial\mathbf{\bar{z}}} & \frac{\partial\Phi_{KL}\left(  \mathbf{z}%
_{0},\mathbf{\bar{z}}_{0}\right)  }{\partial\mathbf{z}}%
\end{array}
\right)  \left(
\begin{array}
[c]{c}%
\delta\mathbf{\bar{z}}\left(  t\right) \\
\delta\mathbf{z}\left(  t\right)
\end{array}
\right)
\end{align}
which have Laplace transforms%
\begin{equation}
\left(
\begin{array}
[c]{cc}%
\frac{\partial\Phi_{KL}\left(  \mathbf{z}_{0},\mathbf{\bar{z}}_{0}\right)
}{\partial\mathbf{\bar{z}}} & \frac{\partial\Phi_{KL}\left(  \mathbf{z}%
_{0},\mathbf{\bar{z}}_{0}\right)  }{\partial\mathbf{z}}%
\end{array}
\right)  \left\{  \left(
\begin{array}
[c]{cc}%
\mathbf{0} & -\xi\mathbf{C}\\
\bar{\xi}\mathbf{\bar{C}} & \mathbf{0}%
\end{array}
\right)  -D^{2}E\left(  \mathbf{z}_{0},\mathbf{\bar{z}}_{0},t\right)
\right\}  ^{-1}D^{1}F\left(  \mathbf{z}_{0},\mathbf{\bar{z}}_{0},\xi\right)
\end{equation}
The matrix elements and the linear response of the 1-particle reduced density
operator can be constructed in the following manner:%
\begin{equation}
D^{N}\left(  \mathbf{z}_{0},\mathbf{\bar{z}}_{0}\right)  =\sum_{K,L}\Phi
_{KL}\left(  \mathbf{z}_{0},\mathbf{\bar{z}}_{0}\right)  \left\vert \Phi
_{L}\right\rangle \left\langle \Phi_{K}\right\vert \Rightarrow D^{1}\left(
\mathbf{z}_{0},\mathbf{\bar{z}}_{0}\right)  _{ij}=\operatorname{Tr}\left\{
D^{N}\left(  \mathbf{z}_{0},\mathbf{\bar{z}}_{0}\right)  a_{i}^{\dagger}%
a_{j}\right\}
\end{equation}
hence
\begin{align}
&  \delta^{1}D^{N}\left(  \mathbf{z}_{0},\mathbf{\bar{z}}_{0}\right)  \left(
t\right)  =\sum_{K,L}\delta^{1}\Phi_{KL}\left(  \mathbf{z}_{0},\mathbf{\bar
{z}}_{0}\right)  \left\vert \Phi_{L}\right\rangle \left\langle \Phi
_{K}\right\vert \Rightarrow\delta^{1}D^{1}\left(  \mathbf{z}_{0}%
,\mathbf{\bar{z}}_{0}\right)  \left(  t\right)  _{ij}=\operatorname{Tr}%
\left\{  \delta D^{N}\left(  \mathbf{z}_{0},\mathbf{\bar{z}}_{0}\right)
\left(  t\right)  a_{i}^{\dagger}a_{j}\right\} \nonumber\\
&  \Rightarrow\delta^{1}D^{1}\left(  \mathbf{z}_{0},\mathbf{\bar{z}}%
_{0}\right)  \left(  t\right)  _{ij}=\operatorname{Tr}\left\{  \sum
_{K,L}\delta^{1}\Phi_{KL}\left(  t\right)  \left(  \mathbf{z}_{0}%
,\mathbf{\bar{z}}_{0}\right)  \left\vert \Phi_{L}\right\rangle \left\langle
\Phi_{K}\right\vert a_{i}^{\dagger}a_{j}\right\} \nonumber\\
&  =\operatorname{Tr}\left\{  \sum_{K,L}\left(
\begin{array}
[c]{cc}%
\frac{\partial\Phi_{KL}\left(  \mathbf{z}_{0},\mathbf{\bar{z}}_{0}\right)
}{\partial\mathbf{\bar{z}}} & \frac{\partial\Phi_{KL}\left(  \mathbf{z}%
_{0},\mathbf{\bar{z}}_{0}\right)  }{\partial\mathbf{z}}%
\end{array}
\right)  \left(
\begin{array}
[c]{c}%
\delta\mathbf{\bar{z}}\left(  t\right) \\
\delta\mathbf{z}\left(  t\right)
\end{array}
\right)  a_{i}^{\dagger}a_{j}\right\} \nonumber\\
&  =\left(
\begin{array}
[c]{cc}%
\frac{\partial D^{1}\left(  \mathbf{z}_{0},\mathbf{\bar{z}}_{0}\right)  _{ij}%
}{\partial\mathbf{\bar{z}}} & \frac{\partial D^{1}\left(  \mathbf{z}%
_{0},\mathbf{\bar{z}}_{0}\right)  _{ij}}{\partial\mathbf{z}}%
\end{array}
\right)  \left(
\begin{array}
[c]{c}%
\delta\mathbf{\bar{z}}\left(  t\right) \\
\delta\mathbf{z}\left(  t\right)
\end{array}
\right)
\end{align}
The Laplace transform of this response is then%
\begin{align}
&  \delta^{1}D^{1}\left(  \mathbf{z}_{0},\mathbf{\bar{z}}_{0}\right)
_{ij}\left(  \xi\right)  =\nonumber\\
&  \left(
\begin{array}
[c]{cc}%
\frac{\partial D^{1}\left(  \mathbf{z}_{0},\mathbf{\bar{z}}_{0}\right)  _{ij}%
}{\partial\mathbf{\bar{z}}} & \frac{\partial D^{1}\left(  \mathbf{z}%
_{0},\mathbf{\bar{z}}_{0}\right)  _{ij}}{\partial\mathbf{z}}%
\end{array}
\right)  \left\{  \left(
\begin{array}
[c]{cc}%
\mathbf{0} & -\xi\mathbf{C}\\
\bar{\xi}\mathbf{\bar{C}} & \mathbf{0}%
\end{array}
\right)  -D^{2}E\left(  \mathbf{z}_{0},\mathbf{\bar{z}}_{0},t\right)
\right\}  ^{-1}D^{1}F\left(  \mathbf{z}_{0},\mathbf{\bar{z}}_{0},\xi\right)
\end{align}


\subsection{Exact Linear Response in Coherent State Setting}

In order to facilitate a direct comparison of approximations to exact linear
response functions it is useful to derive the exact linear response function
in terms of coherent states based on the coset $SU\left(  \mathcal{H}%
^{N}\right)  /U\left(  1\right)  $ The classical equations of motion defined
on this coherent state manifold are equivalent to the full time dependent
Schr\"{o}dinger equation. The configuration space formed by these coherent
sates is constituted by all of the states $\left\{  \left\vert \mathbf{z}%
\right\rangle \right\}  \in\mathcal{H}^{N}$ The group $G=SU\left(
\mathcal{H}^{N}\right)  $ and the subgroup
\begin{equation}
H=\left\{  \left.  g\right\vert g\in G,g\left\vert \Phi_{1}\right\rangle
=e^{i\theta}\left\vert \Phi_{1}\right\rangle \right\}  =\left\{  \exp\left(
i\lambda\left\vert \Phi_{1}\right\rangle \left\langle \Phi_{1}\right\vert
\right)  \exp\left(  i\sum_{i,j\neq1}\lambda_{ij}\left\vert \Phi
_{i}\right\rangle \left\langle \Phi_{j}\right\vert \right)  \right\}
\end{equation}
to produce coset representatives
\begin{equation}
G/H=\left\{  \exp\left(  \sum_{i>1}\eta_{i}\left\vert \Phi_{i}\right\rangle
\left\langle \Phi_{1}\right\vert \right)  \exp\left(  \eta_{1}\left\vert
\Phi_{1}\right\rangle \left\langle \Phi_{1}\right\vert \right)  \exp\left(
-\sum_{i>1}\bar{\eta}_{i}\left\vert \Phi_{i}\right\rangle \left\langle
\Phi_{1}\right\vert \right)  \right\}
\end{equation}
where $\left\{  \left.  \left\vert \Phi_{i}\right\rangle \right\vert
\quad1\leq i\leq\left(
%TCIMACRO{\QATOP{r}{N}}%
%BeginExpansion
\genfrac{}{}{0pt}{}{r}{N}%
%EndExpansion
\right)  \right\}  $ is a complete orthonormal basis for $\mathcal{H}^{N}$.

The family of coherent states associated with this coset is defined as
\begin{equation}
\left\vert \eta\right\rangle =\left\{  \exp\left(  \sum_{i>1}\eta
_{i}\left\vert \Phi_{i}\right\rangle \left\langle \Phi_{1}\right\vert \right)
\exp\left(  \eta_{1}\left\vert \Phi_{i}\right\rangle \left\langle \Phi
_{1}\right\vert \right)  \right\}  \left\vert \Phi_{1}\right\rangle
\end{equation}
for convenience we introduce a new coordinate system $\left\{  \mathbf{z}%
\right\}  $ by setting
\begin{align}
\eta_{1}  &  =\ln\left(  z_{1}\right) \\
\eta_{i}  &  =\frac{z_{i}}{z_{1}};\quad2\leq i\leq\left(
%TCIMACRO{\QATOP{r}{N}}%
%BeginExpansion
\genfrac{}{}{0pt}{}{r}{N}%
%EndExpansion
\right)
\end{align}
which leads to the following form for the coherent states
\begin{equation}
\left\vert \mathbf{z}\right\rangle =\sum_{1\leq K\leq\left(
%TCIMACRO{\QATOP{r}{N}}%
%BeginExpansion
\genfrac{}{}{0pt}{}{r}{N}%
%EndExpansion
\right)  }z_{K}\left\vert \Phi_{K}\right\rangle
\end{equation}
The overlap kernel, classical Hamiltonian and Kaehler metric for this coset
manifold are then defined for an arbitrary quantum Hamiltonian H as:%
\begin{align}
S\left(  \mathbf{z,\bar{z}}\right)   &  =\left\langle \mathbf{z|z}%
\right\rangle =\sum_{1\leq K\leq\left(
%TCIMACRO{\QATOP{r}{N}}%
%BeginExpansion
\genfrac{}{}{0pt}{}{r}{N}%
%EndExpansion
\right)  }\left\vert z_{K}\right\vert ^{2}\\
E\left(  \mathbf{z,\bar{z}}\right)   &  =\frac{\left\langle \mathbf{z|}%
H\mathbf{z}\right\rangle }{\left\langle \mathbf{z|z}\right\rangle }=\frac
{\sum_{1\leq K,L\leq\left(
%TCIMACRO{\QATOP{r}{N}}%
%BeginExpansion
\genfrac{}{}{0pt}{}{r}{N}%
%EndExpansion
\right)  }z_{K}\bar{z}_{L}\left\langle \Phi_{K}\mathbf{|}H\Phi_{L}%
\right\rangle }{\sum_{1\leq K\leq\left(
%TCIMACRO{\QATOP{r}{N}}%
%BeginExpansion
\genfrac{}{}{0pt}{}{r}{N}%
%EndExpansion
\right)  }\left\vert z_{K}\right\vert ^{2}}\\
C_{KL}  &  =\frac{\partial^{2}\ln S}{\partial z_{K}\partial\bar{z}_{L}}%
=\frac{1}{S}\left\{  \frac{\partial^{2}S}{\partial z_{K}\partial\bar{z}_{L}%
}-\frac{1}{S}\frac{\partial S}{\partial z_{K}}\frac{\partial S}{\partial
\bar{z}_{L}}\right\}  =\frac{1}{S}\left\{  \delta_{KL}-\frac{\bar{z}_{K}z_{L}%
}{S}\right\}
\end{align}
this shows that
\begin{equation}
\mathbf{C}=\frac{1}{S}\left\{  \mathbf{I-P}_{z}\right\}  =\frac{1}%
{S}\mathbf{P}_{z}^{\bot}%
\end{equation}
where $\mathbf{P}_{z}$ is the orthogonal projector onto the vector $\left\vert
\mathbf{z}\right\rangle \in\mathbb{C}^{\binom{r}{N}}$. Noting that
\begin{align}
\left(
\begin{array}
[c]{c}%
\frac{\partial E\left(  \mathbf{z},\mathbf{\bar{z}}\right)  }{\partial
\mathbf{\bar{z}}}\\
\frac{\partial E\left(  \mathbf{z},\mathbf{\bar{z}}\right)  }{\partial
\mathbf{z}}%
\end{array}
\right)   &  =\left(
\begin{array}
[c]{c}%
\frac{1}{S\left(  \mathbf{z},\mathbf{\bar{z}}\right)  }\left\{  \frac{\partial
H\left(  \mathbf{z},\mathbf{\bar{z}}\right)  }{\partial\mathbf{z}}-E\left(
\mathbf{z},\mathbf{\bar{z}}\right)  \frac{\partial S\left(  \mathbf{z}%
,\mathbf{\bar{z}}\right)  }{\partial\mathbf{z}}\right\} \\
\frac{1}{S\left(  \mathbf{z},\mathbf{\bar{z}}\right)  }\left\{  \frac{\partial
H\left(  \mathbf{z},\mathbf{\bar{z}}\right)  }{\partial\mathbf{\bar{z}}%
}-E\left(  \mathbf{z},\mathbf{\bar{z}}\right)  \frac{\partial S\left(
\mathbf{z},\mathbf{\bar{z}}\right)  }{\partial\mathbf{\bar{z}}}\right\}
\end{array}
\right) \\
&  =\frac{1}{S\left(  \mathbf{z},\mathbf{\bar{z}}\right)  }\left(
\begin{array}
[c]{c}%
\mathbb{H}\mathbf{\bullet z-}E\mathbf{z}\\
\mathbb{\bar{H}}\mathbf{\bullet\bar{z}-}E\mathbf{\bar{z}}%
\end{array}
\right)  =\frac{1}{S\left(  \mathbf{z},\mathbf{\bar{z}}\right)  }\left(
\begin{array}
[c]{cc}%
\mathbb{O} & \mathbb{H}\mathbf{-}E\mathbb{I}\\
\mathbb{\bar{H}}\mathbf{-}E\mathbb{I} & \mathbb{O}%
\end{array}
\right)  \left(
\begin{array}
[c]{c}%
\mathbf{\bar{z}}\\
\mathbf{z}%
\end{array}
\right)
\end{align}
the classical equations of motion eq. $\left(  \ref{EOM}\right)  $ give
\begin{align}
&  \left(
\begin{array}
[c]{cc}%
\mathbf{0} & \frac{i}{\hslash}\mathbf{C}\\
-\frac{i}{\hslash}\mathbf{\bar{C}} & \mathbf{0}%
\end{array}
\right)  \left(
\begin{array}
[c]{c}%
\frac{d\mathbf{\bar{z}}}{dt}\\
\frac{d\mathbf{z}}{dt}%
\end{array}
\right)  =\left(
\begin{array}
[c]{c}%
\frac{\partial E}{\partial\mathbf{\bar{z}}}\\
\frac{\partial E}{\partial\mathbf{z}}%
\end{array}
\right) \\
&  \Rightarrow\left(
\begin{array}
[c]{cc}%
\mathbf{0} & \frac{i}{\hslash}\mathbf{P}_{z}^{\bot}\\
-\frac{i}{\hslash}\mathbf{P}_{z}^{\bot} & \mathbf{0}%
\end{array}
\right)  \left(
\begin{array}
[c]{c}%
\frac{d\mathbf{\bar{z}}}{dt}\\
\frac{d\mathbf{z}}{dt}%
\end{array}
\right)  =\left(
\begin{array}
[c]{cc}%
\mathbb{O} & \mathbb{H}\mathbf{-}E\mathbb{I}\\
\mathbb{\bar{H}}\mathbf{-}E\mathbb{I} & \mathbb{O}%
\end{array}
\right)  \left(
\begin{array}
[c]{c}%
\mathbf{\bar{z}}\\
\mathbf{z}%
\end{array}
\right)
\end{align}
as $\left(  \mathbb{H}\mathbf{-}E\mathbb{I}\right)  \mathbf{z}$ is also $\bot$
$\mathbf{z}$ this%
\begin{equation}
\left(
\begin{array}
[c]{c}%
-\frac{i}{\hslash}\frac{d\mathbf{\bar{z}}}{dt}\\
\frac{i}{\hslash}\frac{d\mathbf{z}}{dt}%
\end{array}
\right)  =\left(
\begin{array}
[c]{cc}%
\mathbb{O} & \mathbb{I}\\
-\mathbb{I} & \mathbb{O}%
\end{array}
\right)  \left(
\begin{array}
[c]{cc}%
\mathbb{O} & \mathbb{H}\mathbf{-}E\mathbb{I}\\
\mathbb{\bar{H}}\mathbf{-}E\mathbb{I} & \mathbb{O}%
\end{array}
\right)  \left(
\begin{array}
[c]{c}%
\mathbf{\bar{z}}\\
\mathbf{z}%
\end{array}
\right)  =\left(
\begin{array}
[c]{cc}%
\mathbb{\bar{H}}\mathbf{-}E\mathbb{I} & \mathbb{O}\\
\mathbb{O} & \mathbb{H}\mathbf{-}E\mathbb{I}%
\end{array}
\right)  \left(
\begin{array}
[c]{c}%
\mathbf{\bar{z}}\\
\mathbf{z}%
\end{array}
\right)
\end{equation}
the critical point of $E\left(  \mathbf{z},\mathbf{\bar{z}}\right)  $
correspond to the stationary states of $H$ as
\begin{equation}
\Rightarrow\left(
\begin{array}
[c]{c}%
\frac{\partial E}{\partial\mathbf{\bar{z}}}\\
\frac{\partial E}{\partial\mathbf{z}}%
\end{array}
\right)  =\left(
\begin{array}
[c]{c}%
\mathbf{0}\\
\mathbf{0}%
\end{array}
\right)  =\frac{1}{S\left(  \mathbf{z},\mathbf{\bar{z}}\right)  }\left(
\begin{array}
[c]{cc}%
\mathbb{O} & \mathbb{H}\mathbf{-}E\mathbb{I}\\
\mathbb{\bar{H}}\mathbf{-}E\mathbb{I} & \mathbb{O}%
\end{array}
\right)  \left(
\begin{array}
[c]{c}%
\mathbf{\bar{z}}\\
\mathbf{z}%
\end{array}
\right)
\end{equation}
and hence that $\left\vert \mathbf{z}\right\rangle $ are eigen states of $H$.

We now consider the Hamiltonian of the molecular system interacting with a
radiation field. If the point $\left(  \mathbf{z}_{0},\mathbf{\bar{z}}%
_{0}\right)  $ is a critical point of the molecular system i.e. $\left(
\mathbb{H}\mathbf{-}E\mathbb{I}\right)  \bullet\mathbf{z}_{0}=\mathbf{0}$ the
Laplace transform of the linear response function becomes%
\begin{equation}
R\left(  \xi\right)  =\left\{  \left(
\begin{array}
[c]{cc}%
\mathbb{O} & -\xi\mathbf{C}\\
\bar{\xi}\mathbf{\bar{C}} & \mathbb{O}%
\end{array}
\right)  -D^{2}E\left(  \mathbf{z}_{0},\mathbf{\bar{z}}_{0}\right)  \right\}
^{-1}=\left(
\begin{array}
[c]{cc}%
\mathbb{O} & -\left(  \xi\mathbb{I+\bar{H}}\mathbf{-}E\mathbb{I}\right)
^{-1}\\
\left(  \xi\mathbb{I-\bar{H}}\mathbf{-}E\mathbb{I}\right)  ^{-1} & \mathbb{O}%
\end{array}
\right)
\end{equation}
where it \emph{should be observed }that in this case the Hessian matrix is now
block skew diagonal%
\begin{equation}
D^{2}E\left(  \mathbf{z}_{0},\mathbf{\bar{z}}_{0}\right)  =\left(
\begin{array}
[c]{cc}%
\frac{\partial^{2}E\left(  \mathbf{z}_{0},\mathbf{\bar{z}}_{0},t\right)
}{\partial\mathbf{\bar{z}}\partial\mathbf{\bar{z}}} & \frac{\partial
^{2}E\left(  \mathbf{z}_{0},\mathbf{\bar{z}}_{0},t\right)  }{\partial
\mathbf{\bar{z}}\partial\mathbf{z}}\\
\frac{\partial^{2}E\left(  \mathbf{z}_{0},\mathbf{\bar{z}}_{0},t\right)
}{\partial\mathbf{z}\partial\mathbf{\bar{z}}} & \frac{\partial^{2}E\left(
\mathbf{z}_{0},\mathbf{\bar{z}}_{0},t\right)  }{\partial\mathbf{z}%
\partial\mathbf{z}}%
\end{array}
\right)  =\left(
\begin{array}
[c]{cc}%
\mathbb{O} & \frac{\partial^{2}E\left(  \mathbf{z}_{0},\mathbf{\bar{z}}%
_{0},t\right)  }{\partial\mathbf{\bar{z}}\partial\mathbf{z}}\\
\frac{\partial^{2}E\left(  \mathbf{z}_{0},\mathbf{\bar{z}}_{0},t\right)
}{\partial\mathbf{z}\partial\mathbf{\bar{z}}} & \mathbb{O}%
\end{array}
\right)
\end{equation}
The linear response of the FORDM in this case is thus%
\begin{align}
&  \delta D^{1}\left(  \mathbf{z}_{0},\mathbf{\bar{z}}_{0}\right)
_{ij}\left(  \xi\right)  =\\
&  \left(
\begin{array}
[c]{cc}%
\frac{\partial D^{1}\left(  \mathbf{z}_{0},\mathbf{\bar{z}}_{0}\right)  _{ij}%
}{\partial\mathbf{\bar{z}}} & \frac{\partial D^{1}\left(  \mathbf{z}%
_{0},\mathbf{\bar{z}}_{0}\right)  _{ij}}{\partial\mathbf{z}}%
\end{array}
\right)  \left(
\begin{array}
[c]{cc}%
\mathbb{O} & -\left(  \xi\mathbb{I+\bar{H}}\mathbf{-}E\mathbb{I}\right)
^{-1}\\
\left(  \xi\mathbb{I-\bar{H}}\mathbf{+}E\mathbb{I}\right)  ^{-1} & \mathbb{O}%
\end{array}
\right)  D^{1}F\left(  \mathbf{z}_{0},\mathbf{\bar{z}}_{0},\xi\right)
\end{align}
where
\begin{equation}
D^{1}F\left(  \mathbf{z}_{0},\mathbf{\bar{z}}_{0},\xi\right)  =\left(
\begin{array}
[c]{c}%
F\left(  \xi\right)  \mathbf{z}_{0}\\
\bar{F}\left(  \xi\right)  \mathbf{\bar{z}}_{0}%
\end{array}
\right)  =\left(
\begin{array}
[c]{cc}%
\mathbb{O} & \mathbb{F}\left(  \xi\right) \\
\mathbb{\bar{F}}\left(  \xi\right)  \mathbb{I} & \mathbb{O}%
\end{array}
\right)  \left(
\begin{array}
[c]{c}%
\mathbf{\bar{z}}_{0}\\
\mathbf{z}_{0}%
\end{array}
\right)
\end{equation}
The Hamiltonian of the molecular system plus perturbation is
\begin{equation}
H=H_{0}+F\left(  t\right)
\end{equation}
and has the matrix representation wrt to the complete orthonormal basis
$\left\{  \left\vert {\Phi}_{j}\right\rangle \right\}  \equiv\left\{
\left\vert \mathbf{\Phi}\right\rangle \right\}  $%
\begin{equation}
\left\langle \left.  \mathbf{\Phi}\right\vert H\mathbf{\Phi}\right\rangle
=\mathbb{H}=\mathbb{H}_{0}+\mathbb{F}\left(  t\right)
\end{equation}
The derivatives of the energy can be expressed in terms of these matrices as
\begin{align}
\frac{\partial E\left(  \mathbf{z}{_{0},\mathbf{\bar{z}}_{0}}\right)
}{\partial\mathbf{z}}  &  =\frac{1}{S}\left\langle \left.  \mathbf{z}{_{0}%
}\right\vert \left(  H_{0}-E\left(  \mathbf{z}{_{0},\mathbf{\bar{z}}_{0}%
}\right)  I+F\left(  t\right)  \right)  \mathbf{\Phi}\right\rangle \nonumber\\
&  =\frac{1}{S}\left\{  \left\langle \left.  \mathbf{z}{_{0}}\right\vert
\left(  H_{0}-E\left(  \mathbf{z}{_{0},\mathbf{\bar{z}}_{0}}\right)  I\right)
\mathbf{\Phi}\right\rangle +\sum_{kl}f_{kl}{{\left(  t\right)  }}\left\langle
\left.  \mathbf{z}{_{0}}\right\vert a_{k}^{\dag}\left(  t\right)  a_{l}\left(
t\right)  \mathbf{\Phi}\right\rangle \right\} \nonumber\\
&  =\frac{1}{S}\left\{  \mathbb{\bar{H}}_{0}-E\left(  \mathbf{z}%
{_{0},\mathbf{\bar{z}}_{0}}\right)  \mathbb{I}+\mathbb{\bar{F}}\left(
t\right)  \right\}  \bullet\mathbf{z}{_{0}}%
\end{align}
and
\begin{align}
\frac{\partial^{2}E\left(  \mathbf{z}{_{0},\mathbf{\bar{z}}_{0}}\right)
}{\partial\mathbf{z}\partial\mathbf{\bar{z}}}  &  =\frac{1}{S}\left\langle
\left.  \mathbf{\Phi}\right\vert \left(  H_{0}-E\left(  \mathbf{z}%
{_{0},\mathbf{\bar{z}}_{0}}\right)  I+F\left(  t\right)  \right)
\mathbf{\Phi}\right\rangle \nonumber\\
&  =\frac{1}{S}\left\{  \left\langle \left.  \mathbf{\Phi}\right\vert \left(
H_{0}-E\left(  \mathbf{z}{_{0},\mathbf{\bar{z}}_{0}}\right)  I\right)
\mathbf{\Phi}\right\rangle +\sum_{kl}f_{kl}{{\left(  t\right)  }}\left\langle
\left.  \mathbf{\Phi}\right\vert a_{k}^{\dag}\left(  t\right)  a_{l}\left(
t\right)  \mathbf{\Phi}\right\rangle \right\} \nonumber\\
&  =\frac{1}{S}\left\{  \mathbb{\bar{H}}_{0}-E\left(  \mathbf{z}%
{_{0},\mathbf{\bar{z}}_{0}}\right)  \mathbb{I}+\mathbb{\bar{F}}\left(
t\right)  \right\}
\end{align}
Expanding the energy about the point $\left(  \mathbf{z}{_{0},\mathbf{\bar{z}%
}_{0}}\right)  $ gives the \emph{exact} expression
\begin{align}
E\left(  \mathbf{z}{_{0}+\delta\mathbf{z},\mathbf{\bar{z}}_{0}+\delta
\mathbf{\bar{z}}}\right)   &  =E\left(  \mathbf{z}{_{0},\mathbf{\bar{z}}_{0}%
}\right)  +{\delta\mathbf{z}}^{\dagger}\frac{\partial E\left(  \mathbf{z}%
{_{0},\mathbf{\bar{z}}_{0}}\right)  }{\partial{\mathbf{\bar{z}}}}%
+\frac{\partial E\left(  \mathbf{z}{_{0},\mathbf{\bar{z}}_{0}}\right)
}{\partial{\mathbf{z}}}{\delta\mathbf{z}}\nonumber\\
&  +\frac{1}{2}\left\{  {\delta\mathbf{z}}^{\dagger}\frac{\partial^{2}E\left(
\mathbf{z}{_{0},\mathbf{\bar{z}}_{0}}\right)  }{\partial{\mathbf{\bar{z}}%
}\partial{\mathbf{z}}}{\delta\mathbf{z+}\delta\mathbf{\bar{z}}}^{\dagger}%
\frac{\partial^{2}E\left(  \mathbf{z}{_{0},\mathbf{\bar{z}}_{0}}\right)
}{\partial{\mathbf{\bar{z}}}\partial{\mathbf{z}}}{\delta\mathbf{\bar{z}}%
}\right\}
\end{align}
Linear response to the external perturbation is obtained by approximating the
energy function around the point $\left(  \mathbf{z}{_{0},\mathbf{\bar{z}}%
_{0}}\right)  $ by only retaining terms up to first order in the perturbation
\begin{align}
&  E\left(  \mathbf{z}{_{0}+\delta\mathbf{z},\mathbf{\bar{z}}_{0}%
+\delta\mathbf{\bar{z}}}\right)  \approx E_{2}\left(  \mathbf{z}{_{0}%
+\delta\mathbf{z},\mathbf{\bar{z}}_{0}+\delta\mathbf{\bar{z}}}\right)
\nonumber\\
&  ={\delta\mathbf{z}}^{\dagger}\left\langle \left.  \mathbf{\Phi}\right\vert
\left(  H_{0}+F\left(  t\right)  \right)  \mathbf{z}{_{0}}\right\rangle
+\left\langle \left.  \mathbf{z}{_{0}}\right\vert \left(  H_{0}+F\left(
t\right)  \right)  \mathbf{\Phi}\right\rangle {\delta\mathbf{z+}%
\delta\mathbf{z}}^{\dagger}\left\langle \left.  \mathbf{\Phi}\right\vert
H_{0}\mathbf{\Phi}\right\rangle {\delta\mathbf{z}}%
\end{align}
to give the local response energy gradient
\begin{align}
\left(
\begin{array}
[c]{c}%
\frac{\partial E_{2}\left[  \mathbf{z}_{0}\right]  \left(  \delta
\mathbf{z},\delta\mathbf{\bar{z}},t\right)  }{\partial\delta\mathbf{\bar{z}}%
}\\
\frac{\partial E_{2}\left[  \mathbf{z}_{0}\right]  \left(  \delta
\mathbf{z},\delta\mathbf{\bar{z}},t\right)  }{\partial\delta\mathbf{z}}%
\end{array}
\right)   &  =\frac{1}{S}\left(
\begin{array}
[c]{cc}%
\mathbb{H}_{0}-E\mathbb{I+F}\left(  t\right)  & \mathbb{O}\\
\mathbb{O} & \mathbb{\bar{H}}_{0}-E\mathbb{I+\bar{F}}\left(  t\right)
\end{array}
\right)  \left(
\begin{array}
[c]{c}%
\mathbf{z}{_{0}}\\
\mathbf{\bar{z}}{_{0}}%
\end{array}
\right) \nonumber\\
&  +\frac{1}{S}\left(
\begin{array}
[c]{cc}%
\mathbb{H}_{0}-E\mathbb{I} & \mathbb{O}\\
\mathbb{O} & \mathbb{\bar{H}}_{0}-E\mathbb{I}%
\end{array}
\right)  \left(
\begin{array}
[c]{c}%
\delta\mathbf{z}{_{0}}\\
\delta\mathbf{\bar{z}}{_{0}}%
\end{array}
\right)
\end{align}
This implies that
\begin{align}
\delta^{1}D^{1}\left(  \mathbf{z}{_{0},\mathbf{\bar{z}}_{0}}\right)
_{ij}\left(  \xi\right)   &  =\left(  {\mathbf{\bar{z}}_{0}^{\dagger}\cdot
}\overline{\left\langle \left.  \mathbf{\Phi}\right\vert a_{j}^{\dag}%
a_{i}\mathbf{\Phi}\right\rangle },{\mathbf{z}_{0}^{\dagger}\cdot}\left\langle
\left.  \mathbf{\Phi}\right\vert a_{j}^{\dag}a_{i}\mathbf{\Phi}\right\rangle
\right)  \left(
\begin{array}
[c]{cc}%
\mathbb{O} & \mathbb{I}\\
\mathbb{I} & \mathbb{O}%
\end{array}
\right)  \left(
\begin{array}
[c]{cc}%
\mathbb{O} & \mathbb{I}\\
\mathbb{I} & \mathbb{O}%
\end{array}
\right)  \times\nonumber\\
&  \left(
\begin{array}
[c]{cc}%
\mathbb{O} & -\left(  \xi\mathbb{I+\bar{H}}\mathbf{-}E\mathbb{I}\right)
^{-1}\\
\left(  \xi\mathbb{I-\bar{H}}\mathbf{+}E\mathbb{I}\right)  ^{-1} & \mathbb{O}%
\end{array}
\right)  \left(
\begin{array}
[c]{cc}%
\mathbb{F}\left(  \xi\right)  & \mathbb{O}\\
\mathbb{O} & \mathbb{\bar{F}}\left(  \xi\right)  \mathbb{I}%
\end{array}
\right)  \left(
\begin{array}
[c]{c}%
\mathbf{\bar{z}}_{0}\\
\mathbf{z}_{0}%
\end{array}
\right)
\end{align}
which can be expressed as
\begin{align}
&  \left(  \left\langle \left.  {\mathbf{z}_{0}}\right\vert a_{j}^{\dag}%
a_{i}\mathbf{\Phi}\right\rangle \quad\overline{\left\langle \left.
{\mathbf{z}_{0}}\right\vert a_{j}^{\dag}a_{i}\mathbf{\Phi}\right\rangle
}\right)  \times\\
&  \left(
\begin{array}
[c]{cc}%
\left\langle \mathbf{\Phi}\left\vert \left(  \xi I\mathbb{-}H\mathbf{-}%
EI\right)  ^{-1}\right.  \mathbf{\Phi}\right\rangle  & \mathbb{O}\\
\mathbb{O} & \overline{\left\langle \mathbf{\Phi}\left\vert \left(  \xi
I\mathbb{-}H\mathbf{-}EI\right)  ^{-1}\right.  \mathbf{\Phi}\right\rangle }%
\end{array}
\right)  \left(
\begin{array}
[c]{c}%
\left\langle \mathbf{\Phi}\left\vert \left(  \mathbf{a}^{\dagger}%
\mathbf{a}\right)  \right.  {\mathbf{z}_{0}}\right\rangle \\
\overline{\left\langle \mathbf{\Phi}\left\vert \left(  \mathbf{a}^{\dagger
}\mathbf{a}\right)  \right.  {\mathbf{z}_{0}}\right\rangle }%
\end{array}
\right)  \left(
\begin{array}
[c]{c}%
\mathbf{f}\left(  \xi\right) \\
\overline{\mathbf{f}\left(  \xi\right)  }%
\end{array}
\right)  =\\
&  \sum_{k,l}\left\{  \left\langle \left.  {\mathbf{z}_{0}}\right\vert
a_{j}^{\dag}a_{i}\left\{  \left(  \xi I\mathbb{-}H\mathbf{-}EI\right)
^{-1}\right\}  a_{k}^{\dag}a_{l}{\mathbf{z}_{0}}\right\rangle f_{kl}\left(
\xi\right)  -\left\langle \left.  {\mathbf{z}_{0}}\right\vert a_{k}^{\dag
}a_{l}\left\{  \left(  \xi I\mathbb{-}H\mathbf{-}EI\right)  ^{-1}\right\}
a_{j}^{\dag}a_{i}{\mathbf{z}_{0}}\right\rangle \overline{f_{kl}\left(
\xi\right)  }\right\} \\
&  =\sum_{k,l}\left\{  R_{ijkl}\left(  \xi\right)  f_{kl}\left(  \xi\right)
+\overline{R_{ijkl}\left(  \xi\right)  f_{kl}\left(  \xi\right)  }\right\}
\end{align}
Where $R_{ijkl}\left(  \xi\right)  $ are precisely the exact linear response
functions obtained previously.

\subsection{Linear Response Described by $SU\left(  r\right)  /USp\left(
r\right)  $ Coherent States .}

We now use the manifold of AGP coherent states in which to perform the
stationary phase approximation and to consider approximate linear response
functions. Within the context of the classical equations of motion in the
associated phase the linear response of the first order reduced density matrix
produced by the radiation field is then given by
\begin{align}
&  \delta^{1}D^{1}\left(  \mathbf{z}{_{0},\mathbf{\bar{z}}_{0}}\right)
_{ij}\left(  \xi\right)  =\left(
\begin{array}
[c]{cc}%
\frac{\partial D^{1}\left(  \mathbf{z}_{0},\mathbf{\bar{z}}_{0}\right)  _{ij}%
}{\partial\mathbf{\bar{z}}} & \frac{\partial D^{1}\left(  \mathbf{z}%
_{0},\mathbf{\bar{z}}_{0}\right)  _{ij}}{\partial\mathbf{z}}%
\end{array}
\right)  \left(
\begin{array}
[c]{cc}%
\mathbb{O} & \mathbb{I}\\
\mathbb{I} & \mathbb{O}%
\end{array}
\right)  \left(
\begin{array}
[c]{cc}%
\mathbb{O} & \mathbb{I}\\
\mathbb{I} & \mathbb{O}%
\end{array}
\right)  \times\nonumber\\
&  \quad\quad\quad\left(
\begin{array}
[c]{cc}%
\mathbb{O} & \mathbb{I}\\
\mathbb{I} & \mathbb{O}%
\end{array}
\right)  \left(
\begin{array}
[c]{cc}%
\mathbb{O} & \mathbb{I}\\
\mathbb{I} & \mathbb{O}%
\end{array}
\right)  \left\{  \left(
\begin{array}
[c]{cc}%
\mathbb{O} & -\xi\mathbf{C}\\
\bar{\xi}\mathbf{\bar{C}} & \mathbb{O}%
\end{array}
\right)  -D^{2}E\left(  \mathbf{z}_{0},\mathbf{\bar{z}}_{0}\right)  \right\}
^{-1}D^{1}F\left(  \mathbf{z}_{0},\mathbf{\bar{z}}_{0},\xi\right) \nonumber\\
&  =\left(
\begin{array}
[c]{cc}%
\frac{\partial D^{1}\left(  \mathbf{z}_{0},\mathbf{\bar{z}}_{0}\right)  _{ij}%
}{\partial\mathbf{\bar{z}}} & \frac{\partial D^{1}\left(  \mathbf{z}%
_{0},\mathbf{\bar{z}}_{0}\right)  _{ij}}{\partial\mathbf{z}}%
\end{array}
\right)  \times\nonumber\\
&  \quad\quad\quad\quad\quad\quad\left\{  \left(
\begin{array}
[c]{cc}%
\bar{\xi}\mathbf{\bar{C}} & \mathbb{O}\\
\mathbb{O} & -\xi\mathbf{C}%
\end{array}
\right)  -\left(
\begin{array}
[c]{cc}%
\frac{\partial^{2}E\left(  \mathbf{z}_{0},\mathbf{\bar{z}}_{0},t\right)
}{\partial\mathbf{\bar{z}}\partial\mathbf{\bar{z}}} & \frac{\partial
^{2}E\left(  \mathbf{z}_{0},\mathbf{\bar{z}}_{0},t\right)  }{\partial
\mathbf{\bar{z}}\partial\mathbf{z}}\\
\frac{\partial^{2}E\left(  \mathbf{z}_{0},\mathbf{\bar{z}}_{0},t\right)
}{\partial\mathbf{z}\partial\mathbf{\bar{z}}} & \frac{\partial^{2}E\left(
\mathbf{z}_{0},\mathbf{\bar{z}}_{0},t\right)  }{\partial\mathbf{z}%
\partial\mathbf{z}}%
\end{array}
\right)  \right\}  ^{-1}\left(
\begin{array}
[c]{c}%
\frac{\partial F\left(  \mathbf{z},\mathbf{\bar{z},\xi}\right)  }%
{\partial\mathbf{\bar{z}}}\\
\frac{\partial F\left(  \mathbf{z},\mathbf{\bar{z},\xi}\right)  }%
{\partial\mathbf{z}}%
\end{array}
\right)
\end{align}
as
\begin{align}
\frac{\partial F\left(  \mathbf{z}_{0},\mathbf{\bar{z}}_{0}\mathbf{,\xi
}\right)  }{\partial\mathbf{z}}  &  =\sum_{k,l}\left.  \frac{\partial
\left\langle \left.  \mathbf{z}\right\vert a_{k}^{\dag}a_{l}\mathbf{z}%
\right\rangle }{\partial\mathbf{z}}\right\vert _{\mathbf{z=z}_{0}}%
f_{kl}\left(  \xi\right) \\
\frac{\partial D^{1}\left(  \mathbf{z}_{0},\mathbf{\bar{z}}_{0}\right)  _{lk}%
}{\partial\mathbf{z}}  &  =\left.  \frac{\partial\left\langle \left.
\mathbf{z}\right\vert a_{k}^{\dag}a_{l}\mathbf{z}\right\rangle }%
{\partial\mathbf{z}}\right\vert _{\mathbf{z=z}_{0}}%
\end{align}
we obtain
\begin{align}
&  \delta^{1}D^{1}\left(  \mathbf{z}{_{0},\mathbf{\bar{z}}_{0}}\right)
_{ij}\left(  \xi\right)  =\left(
\begin{array}
[c]{cc}%
\left.  \frac{\partial\left\langle \left.  \mathbf{z}\right\vert a_{j}^{\dag
}a_{i}\mathbf{z}\right\rangle }{\partial\mathbf{\bar{z}}}\right\vert
_{\mathbf{z=z}_{0}} & \left.  \frac{\partial\left\langle \left.
\mathbf{z}\right\vert a_{j}^{\dag}a_{i}\mathbf{z}\right\rangle }%
{\partial\mathbf{z}}\right\vert _{\mathbf{z=z}_{0}}%
\end{array}
\right)  \times\nonumber\\
&  \quad\quad\quad\quad\quad\quad\left\{  \left(
\begin{array}
[c]{cc}%
\bar{\xi}\mathbf{\bar{C}} & \mathbb{O}\\
\mathbb{O} & -\xi\mathbf{C}%
\end{array}
\right)  -\left(
\begin{array}
[c]{cc}%
\frac{\partial^{2}E\left(  \mathbf{z}_{0},\mathbf{\bar{z}}_{0},t\right)
}{\partial\mathbf{\bar{z}}\partial\mathbf{z}} & \frac{\partial^{2}E\left(
\mathbf{z}_{0},\mathbf{\bar{z}}_{0},t\right)  }{\partial\mathbf{\bar{z}%
}\partial\mathbf{\bar{z}}}\\
\frac{\partial^{2}E\left(  \mathbf{z}_{0},\mathbf{\bar{z}}_{0},t\right)
}{\partial\mathbf{z}\partial\mathbf{z}} & \frac{\partial^{2}E\left(
\mathbf{z}_{0},\mathbf{\bar{z}}_{0},t\right)  }{\partial\mathbf{z}%
\partial\mathbf{\bar{z}}}%
\end{array}
\right)  \right\}  ^{-1}\times\nonumber\\
&  \quad\quad\quad\quad\quad\quad\quad\quad\quad\quad\quad\quad\left(
\begin{array}
[c]{cc}%
\left.  \frac{\partial\left\langle \left.  \mathbf{z}\right\vert
\mathbf{a}^{\dag}\mathbf{az}\right\rangle }{\partial\mathbf{z}}\right\vert
_{\mathbf{z=z}_{0}} & \mathbb{O}\\
\mathbb{O} & \left.  \frac{\partial\left\langle \left.  \mathbf{z}\right\vert
\mathbf{a}^{\dag}\mathbf{az}\right\rangle }{\partial\mathbf{\bar{z}}%
}\right\vert _{\mathbf{z=z}_{0}}%
\end{array}
\right)  \left(
\begin{array}
[c]{c}%
\mathbf{f}\left(  \xi\right) \\
\overline{\mathbf{f}\left(  \xi\right)  }%
\end{array}
\right)
\end{align}
In general the off diagonal blocks $\left\{  \frac{\partial^{2}E\left(
\mathbf{z}_{0},\mathbf{\bar{z}}_{0},t\right)  }{\partial\mathbf{\bar{z}%
}\partial\mathbf{\bar{z}}},\frac{\partial^{2}E\left(  \mathbf{z}%
_{0},\mathbf{\bar{z}}_{0},t\right)  }{\partial\mathbf{z}\partial\mathbf{z}%
}\right\}  $ are not zero. This can be seen by considering
\begin{align}
H_{0}\left(  \mathbf{z}{,\mathbf{\bar{z}}}\right)   &  =\left\langle \left.
\exp\left\{  \sum\limits_{_{\substack{-s\leq j<k\leq s\\\left\vert
j\right\vert \neq\left\vert k\right\vert }}}z_{+jk}q_{jk}^{\dagger}\right\}
\exp\left\{  \sum\limits_{_{\substack{-s\leq j<k\leq s\\\left\vert
j\right\vert \neq\left\vert k\right\vert }}}\lambda_{jk}u_{jk}\right\}
g_{0}^{\frac{N}{2}}\right\vert H_{0}\right. \nonumber\\
&  \quad\quad\quad\quad\quad\quad\quad\quad\quad\quad\quad\quad\left\vert
\exp\left\{  \sum\limits_{_{\substack{-s\leq j<k\leq s\\\left\vert
j\right\vert \neq\left\vert k\right\vert }}}z_{+jk}q_{jk}^{\dagger}\right\}
\exp\left\{  \sum\limits_{_{\substack{-s\leq j<k\leq s\\\left\vert
j\right\vert \neq\left\vert k\right\vert }}}\lambda_{jk}u_{jk}\right\}
g_{0}^{\frac{N}{2}}\right\rangle
\end{align}
one can see that
\begin{equation}
\frac{\partial^{2}H_{0}\left(  \mathbf{z}{,\mathbf{\bar{z}}}\right)
}{\partial\mathbf{z}_{+}\partial\mathbf{z}_{+}}=0\Rightarrow\left\langle
\left.  g_{0}^{\frac{N}{2}}\left(  \mathbf{z}_{0}\right)  \right\vert
H_{0}q_{i_{1}j_{1}}^{\dagger}\cdots q_{i_{m}j_{m}}^{\dagger}g_{0}^{\frac{N}%
{2}}\left(  \mathbf{z}_{0}\right)  \right\rangle =0\quad\forall m\geq2
\end{equation}
implies that $H_{0}\left(  \mathbf{z}{,\mathbf{\bar{z}}}\right)  $ is first
order in $\mathbf{z}$ and that $H_{0}$ is a linear function of $\left\{
q_{i_{1}j_{1}}^{\dagger}\right\}  $. From
\begin{equation}
\frac{\partial^{2}H_{0}\left(  \mathbf{z}{,\mathbf{\bar{z}}}\right)
}{\partial\mathbf{\bar{z}}_{+}\partial\mathbf{\bar{z}}_{+}}=0,\frac
{\partial^{2}H_{0}\left(  \mathbf{z}{,\mathbf{\bar{z}}}\right)  }%
{\partial\mathbf{\lambda}\partial\mathbf{\lambda}}=0,\frac{\partial^{2}%
H_{0}\left(  \mathbf{z}{,\mathbf{\bar{z}}}\right)  }{\partial\mathbf{\bar
{\lambda}}\partial\mathbf{\bar{\lambda}}}=0
\end{equation}
one reaches the same conclusion about the operators $\left\{  q_{ij}\right\}
$ and $\left\{  \lambda_{ij}\right\}  .$ If this were true then
\begin{equation}
H_{0}=E_{0}\left\vert g_{0}^{\frac{N}{2}}\right\rangle \left\langle
g_{0}^{\frac{N}{2}}\right\vert +\sum_{n}\left(  E_{n}-E_{0}\right)
Q_{n}^{\dagger}\left\vert g_{0}^{\frac{N}{2}}\right\rangle \left\langle
g_{0}^{\frac{N}{2}}\right\vert Q_{n}%
\end{equation}
where
\begin{equation}
Q_{n}^{\dagger}=\sum\limits_{_{\substack{-s\leq j<k\leq s\\\left\vert
j\right\vert \neq\left\vert k\right\vert }}}Q_{njk}q_{jk}^{\dagger}%
+\sum\limits_{_{\substack{-s\leq j<k\leq s\\\left\vert j\right\vert
\neq\left\vert k\right\vert }}}U_{njk}u_{jk}%
\end{equation}
Thus if the excited states were single excitations out of the reference state
were exact, the off diagonal blocks would be identically zero and the
Hamiltonian could be written as a bilinear form in the above excitation
operators. In this case the quadratic expansion would be exact and Taylor
series quadratic error would be zero. (Note that diagonalizing the Hessian
does not accomplish this unless the off diagonal blocks are actually zero, as
the eigen excitation operators would then be linear combinations of the
excitation operators and their adjoints.) The slower the convergence of the
Taylor series, i.e. the greater the quadratic error term, the poorer the
approximation of excited states by single particle excitations out of the
reference is.

Once the coherent state manifold is defined one is free to use any well
defined coordinate system in which to find the critical points of the energy
function i.e. classical Hamiltonian, for that manifold. We have published a
procedure based on the coordinates based on the following parameterization of
AGP states%
\begin{equation}
\left\vert g^{\frac{N}{2}}\right\rangle =P_{N}\exp\left\{  \sum_{1\leq i<j\leq
r}g_{ij}a_{i}^{\dagger}a_{j}^{\dagger}\right\}  \left\vert \phi\right\rangle
\end{equation}
This procedure will be compared to ones based on coordinate system \{z\}.

\section{Appendix}

\section{Lie Algebras\label{liealg}}

We will base all expressions on the standard second quantized formulation
electronic quantum mechanics, i.e. Fermi Fock space which is the direct sum
\begin{equation}
\mathcal{H}_{F}=%
%TCIMACRO{\tbigoplus \limits_{0\leq N\leq r}}%
%BeginExpansion
{\textstyle\bigoplus\limits_{0\leq N\leq r}}
%EndExpansion
\mathcal{H}^{N}%
\end{equation}
where
\begin{equation}
\mathcal{H}^{N}=%
%TCIMACRO{\tbigwedge }%
%BeginExpansion
{\textstyle\bigwedge}
%EndExpansion
\mathcal{H}^{1}%
\end{equation}
is the $N$ -fold antisymmetric tensor product of the one particle state space
$\mathcal{H}^{1}$ which is generated by the basis of one particle functions of
position and spin
\begin{equation}
\left\{  \left.  \varphi_{j}\right\vert 1\leq j\leq r\right\}
\end{equation}
The Banach space of bounded operators acting in $\mathcal{H}_{F}$ are denoted
by $\mathcal{B}\left(  \mathcal{H}_{F}\right)  $ and this space has a second
quantized operator basis
\begin{equation}
\left\{  \left.  a_{j_{1}}^{\dagger}\cdots a_{j_{p}}^{\dagger}a_{k_{q}}\cdots
a_{k_{1}}\right\vert 1\leq j_{1}<\cdots j_{p}\leq r;1\leq k_{1}<\cdots
k_{p}\leq r;1\leq p,q\leq r\right\}
\end{equation}
where
\begin{align}
a_{j}  &  =a\left(  \varphi_{j}\right) \\
a_{j}^{\dagger}  &  =a^{\dagger}\left(  \varphi_{j}\right)  =a\left(
\varphi_{j}\right)  ^{\dagger}%
\end{align}
that whose generators satisfy the Canonical Anti commutation Relationships
(CAR)
\begin{equation}
\left[  a_{j}^{\dagger},a_{k}\right]  _{+}=\left\langle \varphi_{j}%
|\varphi_{k}\right\rangle I.
\end{equation}
The generators
\begin{equation}
\left\{  a_{j}^{\dagger},a_{k},a_{j}^{\dagger}a_{k}-\frac{1}{2}\delta
_{jk}I,a_{j}^{\dagger}a_{k}^{\dagger},a_{j}a_{k};1\leq j\leq k\leq r\right\}
\end{equation}
satisfy the commutation relationships of the lie algebra $so\left(
2r+1\right)  $ and generate a representation of the universal enveloping
algebra $so\left(  2r+1\right)  $ in $\mathcal{B}\left(  \mathcal{H}%
_{F}\right)  ,$ which is fact all of $\mathcal{B}\left(  \mathcal{H}%
_{F}\right)  .$ The subset of generators
\begin{equation}
\left\{  a_{j}^{\dagger}a_{k}-\frac{1}{2}\delta_{jk}I,a_{j}^{\dagger}%
a_{k}^{\dagger},a_{j}a_{k};1\leq j\leq k\leq r\right\}
\end{equation}
satisfy the commutation relationships of the lie algebra of $so\left(
2r\right)  $ and generate a representation of the universal enveloping algebra
$so\left(  2r\right)  $ in $\mathcal{B}\left(  \mathcal{H}_{F}\right)  .$ The
generators
\begin{equation}
\left\{  a_{j}^{\dagger}a_{k}-\frac{1}{2}\delta_{jk}I;1\leq j\leq k\leq
r\right\}
\end{equation}
satisfy the commutation relationships of the lie algebra $su\left(  r\right)
$ and generate a representation of the universal enveloping algebra $su\left(
r\right)  $ in $\mathcal{B}\left(  \mathcal{H}_{F}\right)  $. The
representation of $so\left(  2r+1\right)  $ is irreducible, while the
representations of $so\left(  2r\right)  $ and $su\left(  r\right)  $ are
reducible. The irreducible representations of $so\left(  2r\right)  $ are
obtained by restricting to the subspace of even or odd number states in
$\mathcal{H}_{F},$ while the irreducible representations of $su\left(
r\right)  $ are obtained by restricting to subspaces of states of fixed
integer particle number.


\end{document}